  %%%%%%%%%%%%%%%%%%%%%%%%%%%%%%%%%%%%%%%%%%%%%%%%%%%%%%%%%%%%%%%%%%%%%%%%%%%%
%% Author template for Management Science (mnsc) for articles with e-companion (EC)
%% Mirko Janc, Ph.D., INFORMS, mirko.janc@informs.org
%% ver. 0.95, December 2010
%%%%%%%%%%%%%%%%%%%%%%%%%%%%%%%%%%%%%%%%%%%%%%%%%%%%%%%%%%%%%%%%%%%%%%%%%%%%
%\documentclass[a4paper,UKenglish,cleveref, autoref, thm-restate]{informs3} % current default for manuscript submission
\documentclass[a4paper,UKenglish,cleveref, autoref, thm-restate]{informs3}

\OneAndAHalfSpacedXI % current default line spacing
%%\OneAndAHalfSpacedXII 
%%\DoubleSpacedXII
%%\DoubleSpacedXI

% If hyperref is used, dvi-to-ps driver of choice must be declared as
%   an additional option to the \documentstyle. For example
%\documentclass[dvips,mnsc]{informs3}      % if dvips is used
%\documentclass[dvipsone,mnsc]{informs3}   % if dvipsone is used, etc.

% Private macros here (check that there is no clash with the style)

% Natbib setup for author-year style
\usepackage[authoryear]{natbib}
 \bibpunct[, ]{(}{)}{,}{a}{}{,}%
 \def\bibfont{\small}%
 \def\bibsep{\smallskipamount}%
 \def\bibhang{24pt}%
 \def\newblock{\ }%
 \def\BIBand{and}%

% In preamble
\usepackage[section]{placeins}

\usepackage{float}
\usepackage[utf8]{inputenc}
\usepackage{amssymb}
\usepackage{booktabs}
\usepackage{graphicx}
\usepackage{subcaption}
\usepackage{adjustbox}
\usepackage{multirow}
\usepackage[table,xcdraw]{xcolor}
\usepackage[colorlinks=true, linkcolor=blue, citecolor=blue, urlcolor=blue]{hyperref}
\usepackage{verbatim}
\usepackage{threeparttable}  % add to preamble
\usepackage{booktabs}        % if not already there
\usepackage{enumitem}
\usepackage{caption}
\captionsetup{justification=raggedright, singlelinecheck=false, format=plain}

\usepackage{tabularx}
\newcolumntype{R}{>{\raggedleft\arraybackslash}X}

\newcommand{\julnote}[1]{{\color{purple}{{\bf Julien:} #1}}}
\newcommand{\jnote}[1]{}
% Revision marker: new/changed material shown in purple.
\newcommand{\rev}[1]{{\color{purple}#1}}
% Latest approved targeted revisions: red, without changing earlier purple markings.
\DeclareRobustCommand{\latestrev}[1]{{\color{red}\hypersetup{linkcolor=red,citecolor=red,urlcolor=red}#1}}
\newenvironment{latestrevisions}{\begingroup\color{red}\hypersetup{linkcolor=red,citecolor=red,urlcolor=red}}{\endgroup}
\pdfstringdefDisableCommands{\def\rev#1{#1}\def\latestrev#1{#1}}

\newcommand{\boldpara}[1]{\noindent\textbf{#1}\quad}

\newcommand{\appendix}{%
  \setcounter{section}{0}%
  \renewcommand{\thesection}{Appendix~\Alph{section}}%
  \renewcommand{\theHsection}{Appendix.\Alph{section}}%
}
\renewcommand{\footnoterule}{%
  \kern -3pt
  \hrule width \textwidth height 0.4pt
  \kern 3pt
}

%\bibliographystyle{plainurl}% the mandatory bibstyle
\bibliographystyle{ecta}
%% Setup of theorem styles. Outcomment only one.
%% Preferred default is the first option.
\TheoremsNumberedThrough     % Preferred (Theorem 1, Lemma 1, Theorem 2)
%\TheoremsNumberedByChapter  % (Theorem 1.1, Lema 1.1, Theorem 1.2)
\ECRepeatTheorems

%% Setup of the equation numbering system. Outcomment only one.
%% Preferred default is the first option.
\EquationsNumberedThrough    % Default: (1), (2), ...
%\EquationsNumberedBySection % (1.1), (1.2), ...

% For new submissions, leave this number blank.
% For revisions, input the manuscript number assigned by the on-line
% system along with a suffix ".Rx" where x is the revision number.
%\MANUSCRIPTNO{MS-0001-1922.65}

% The INFORMS banner macros are commented out in this customized informs3.cls
% and only defined under blindrev/nonblindrev; define them empty so \maketitle works.
\providecommand{\theARTICLETOPLEFT}{}
\providecommand{\theARTICLETOPRIGHT}{}

%%%%%%%%%%%%%%%%
\begin{document}
%%%%%%%%%%%%%%%%
\setcounter{page}{2}

% Outcomment only when entries are known. Otherwise leave as is and
%   default values will be used.
%\setcounter{page}{1}
%\VOLUME{00}%
%\NO{0}%
%\MONTH{Xxxxx}% (month or a similar seasonal id)
%\YEAR{0000}% e.g., 2005
%\FIRSTPAGE{000}%
%\LASTPAGE{000}%
%\SHORTYEAR{00}% shortened year (two-digit)
%\ISSUE{0000} %
%\LONGFIRSTPAGE{0001} %
%\DOI{10.1287/xxxx.0000.0000}%

% Author's names for the running heads
% Sample depending on the number of authors;
% \RUNAUTHOR{Jones}
% \RUNAUTHOR{Jones and Wilson}
% \RUNAUTHOR{Jones, Miller, and Wilson}
% \RUNAUTHOR{Jones et al.} % for four or more authors
% Enter authors following the given pattern:
\RUNAUTHOR{Shock Propagation in Decentralized Lending Networks}

% Title or shortened title suitable for running heads. Sample:
% \RUNTITLE{Bundling Information Goods of Decreasing Value}
% Enter the (shortened) title:
\RUNTITLE{Evidence from the Compound Protocol}



\TITLE{Shock Propagation in Decentralized Lending Networks: Evidence from the Compound Protocol
% \thanks{\setlength{\baselineskip}{10pt} We are grateful to Co-Pierre Georg, Wenqian Huang and Riho Marten Pallum
% for their valuable comments and suggestions. 
% We also thank conference participants at the Futures of Payment (Bocconi University), CEPR FinTech and Digital Currency, Future of Money II, Grenoble DeFi (INP-UGA), ICCF 2024, Frontiers in DeFi (Complexity Science Hub), Annual Conference on DeFi \& Digital Currencies (WBS Gillmore Centre), the Economics of Decentralisation (Waseda) for helpful discussions. Julien Prat gratefully acknowledges financial 
% support from the ANR project “BLOCKFI: Blockchain and Decentralized Finance”- 2024 - CES38. The authors acknowledge the use of LLM assistants to support the writing and coding process.}
}

% Block of authors and their affiliations starts here:
% NOTE: Authors with same affiliation, if the order of authors allows,
%   should be entered in ONE field, separated by a comma.
%   \EMAIL field can be repeated if more than one author
% \ARTICLEAUTHORS{%
% \AUTHOR{Natkamon Tovanich}
% \AFF{TU Wien, Austria, \EMAIL{natkamon.tovanich@tuwien.ac.at}}
% \AUTHOR{Stefania Marcassa}
% \AFF{THEMA (UMR CNRS 8184), CY Cergy Paris Université, France, \EMAIL{stefania.marcassa@cyu.fr}} %, \URL{}}
% \AUTHOR{Stefan Kitzler}
% \AFF{Complexity Science Hub, Austria
% and AIT Austrian Institute of Technology, Austria, \EMAIL{kitzler@csh.ac.at}}
% \AUTHOR{Christos A. Makridis}
% \AFF{Arizona State University, AZ, USA; University of Nicosia, Cyprus, \EMAIL{cmakridi@asu.edu}} %, \URL{}}
% \AUTHOR{Julien Prat}
% \AFF{CREST, CNRS, École Polytechnique, Institut Polytechnique de Paris, France, \EMAIL{julien.prat@ensae.fr}}} 

%\authorrunning{Tovanich et al.} %TODO mandatory. First: Use abbreviated first/middle names. Second (only in severe cases): Use first author plus 'et al.'
\ABSTRACT{%
\textbf{Abstract:} We study the propagation of financial shocks in decentralized finance (DeFi) using a novel dataset from Compound, a prominent DeFi lending application on Ethereum. Unlike in traditional interbank networks, Compound's users lend and borrow via liquidity pools implemented as smart contracts. We construct daily balance sheets for users and pools from January 2020 to June 2024, model the liability network, and apply a propagation algorithm that captures the impact of liquidations following token-specific price shocks. We assess our model’s ability to predict financial stress and show \rev{that the indicators published by the protocol explain less than a quarter of the out-of-sample variation in simulated \latestrev{liquidity} stress, whereas measures built from the liability network raise explained variation above 80 percent.} These results underscore the need for \latestrev{network-based} risk monitoring in algorithmic credit systems.
}%

\KEYWORDS{decentralized finance; systemic risk; financial networks; decentralized lending. \textit{JEL}: D85, G01, G23.}

%\textit{JEL}: D85, G01, G23

%\HISTORY{This paper was
%first submitted on June 30, 2025.}

\maketitle

%\input{new_intro.tex}

\section{Introduction}

\rev{A promise of decentralized finance (DeFi) is that it mitigates
some of the risks inherent in financial intermediation. Because positions are
recorded on a public ledger and continuously monitored by smart contracts,
decentralized protocols are held to reduce the counterparty risk
that arises when lenders must rely on trusted
intermediaries.\footnote{\cite{schar2022defi} and \cite{harvey2021defi}
discuss how smart contracts can replace trusted intermediaries and
lower counterparty risk. \cite{aramonte2021defi} offer an early
counterpoint, arguing that leverage, liquidity mismatches and
interconnectedness leave decentralized protocols exposed despite these
features.} \latestrev{These features raise the question of how much risk can be monitored using standard pool-level indicators. We examine this question} by reconstructing the liability network of Compound, one of the
first decentralized lending protocols, and simulating the propagation of price shocks through it. We find that the \latestrev{standard pool-level} indicators published by the protocol predict only a fraction of the \latestrev{simulated} stress. Systemic risk does not vanish under
decentralization: it is relocated to the positions that users hold
across different assets, and \latestrev{monitoring these connections requires information beyond the standard pool-level indicators}.}

\rev{The literature on financial networks has shown that the severity of shock propagation depends critically on the topology of financial linkages.}\footnote{Seminal models of financial contagion include \cite{allen2000financial} and \cite{elliott2014financial}, while \cite{glasserman2016contagion} provide a comprehensive survey of the early literature on the topic. See also \cite{acemoglu2015systemic} and \cite{gabaix2011granular} for theoretical results on the propagation of systemic and idiosyncratic risk in networks, and \cite{digiovanni2014firms} or \cite{carvalho2021supply} for empirical results from micro-data.} \rev{Fragility is determined not only by the leverage and the size of individual institutions, but also by the structure of the network in which they are embedded: a local shock may remain contained when exposures are diffuse and well capitalized, yet propagate into systemic stress when it reaches highly connected counterparties.}

This paper studies how these insights can be applied to supervise and predict systemic risk in the expanding ecosystem of \rev{DeFi}. Unlike traditional financial institutions, DeFi platforms operate through algorithmic governance and smart contracts~\citep{capponi2023defi}, enabling a wide range of financial activities to be conducted without centralized oversight. The novelty of DeFi’s economic design, together with the very features that foster innovation—namely open access, transparency, and composability—also expand the attack surface of these protocols. This increased exposure renders them susceptible to exploits that can precipitate abrupt liquidity crises.\footnote{Among the risks specific to DeFi, the most extensively studied include smart contract exploits~\citep{Sun2021}, flash loan attacks~\citep{gudgeon2020deficrisis}, and governance challenges~\citep{Dotan2023a}.}

Crucially, these risks are no longer confined to the DeFi ecosystem. As regulated financial institutions deepen their integration with decentralized infrastructures through tokenized assets and central bank digital currencies, understanding how systemic vulnerabilities emerge and propagate within digital financial networks becomes increasingly important.

The dynamics of risk propagation in DeFi differ from those in traditional finance: whereas interbank contagion is dampened by balance sheet opacity and legal frictions~\citep{Gennaioli2013ShadowBanking}, risk transmission in DeFi occurs through automated mechanisms that operate almost instantaneously. Such algorithmic propagation reduces the buffering effect of discretionary intervention and may therefore accelerate the spread of financial distress. The specificity and novelty of these amplification channels motivate our analysis of shock propagation within DeFi.

Owing to their market size and inherent risk, lending protocols stand out as the most  natural candidates for the study of financial contagion in DeFi. With total value locked (TVL) exceeding \$30 billion at the time of writing,\footnote{See \url{https://coinmarketcap.com/currencies/compound/} 
and \url{https://www.coingecko.com/en/coins/compound} for historical data on the market 
capitalization of Compound.} lending and borrowing are among the most popular activities for DeFi users. These  activities are also among the riskiest, as borrowers are exposed to liquidation. Among  the various lending protocols, we select Compound for our empirical analysis because it 
has served as a template for numerous protocols, making its architecture relevant beyond our specific context. Furthermore, as one of the earliest decentralized lending platforms, Compound also offers one of the longest available data series~\citep{leshner2019compound}.

In Compound’s financial network, lending pools serve as intermediaries between lenders and borrowers, \latestrev{linking users to pools rather than pools directly to one another}. Instead of lending to one another, users deposit tokens into asset-specific pools and receive tokenized claims, while borrowers draw from those pools after posting collateral. This structure means that distress does not propagate through bilateral pool-to-pool exposures, but through links created by the lending and borrowing activity of users. The mechanism differs from standard interbank contagion. In Compound, distress moves through collateral constraints, automated liquidations, reserve withdrawals, and residual bad debt rather than through direct bilateral default among financial institutions.

Capturing this propagation mechanism requires reconstructing the full network of liabilities. Our first contribution is therefore methodological: we detail a procedure for rebuilding pools' and users' balance sheets and show how these can be used to map the network through which financial stress propagates. This representation allows us to identify and characterize the two main sources of risk resulting from a drop in the market value of tokens used as collateral. First, price shocks induce \textit{liquidity risk} as liquidators deplete the reserve of the shocked pool by \latestrev{seizing and redeeming claims on collateral tokens}. Second, price shocks may generate \textit{solvency risk} for the protocol by creating bad debt when borrowers' positions become undercollateralized. We explain how the design of Compound is largely motivated by the need to contain these two sources of risk. Having clarified the economic logic of the protocol, we propose a propagation algorithm that captures the impact of price shocks, providing a tractable framework for stress testing.

We then apply our methodology to an empirical analysis of Compound. Using the protocol’s indexed subgraph provided by The Graph,\footnote{The Graph is a decentralized indexing protocol that organizes and stores blockchain data through open APIs known as subgraphs. It enables 
developers to query on-chain data from protocols such as Compound in real time, without having to process raw blockchain transactions directly. See \url{https://thegraph.com/} for more information.} we reconstruct daily snapshots of Compound’s financial network from January 1, 2020, to June 15, 2024. The data indicate that the network evolved in ways that altered its vulnerability to shocks. Early periods (2020–2021) are characterized by dense user participation and elevated leverage, whereas later phases (from 2022 onward) exhibit reduced activity. Despite these changes, five pools—ETH, WBTC, USDC, DAI, and USDT—retain a preponderant role due to their substantially larger size relative to other pools. We also document marked differences in the use of cryptoassets versus stablecoins, with the former serving primarily as collateral and the latter accounting for the bulk of borrowing activity.

The specific topology of the financial network, together with the automated nature of liquidations, gives rise to novel risk dynamics that require adapting existing stress-testing frameworks. To this end, we apply our propagation algorithm to historical realizations of Compound’s network in order to generate time series measuring the effects of price shocks on the main tokens. We use these series to construct measures of liquidity and solvency risk, which in turn enable us to evaluate an operational framework aimed at \textit{predicting} these two risk measures. \latestrev{We combine each pool's own indicators with characteristics of its users, weighted by the size of their positions. This identifies which observable summaries of the network retain information about simulated stress.} Leveraging daily snapshots of the lending network, we examine how the financial and topological features of pools and users contribute to the liquidity and solvency risk associated with each token, demonstrating through a series of tests the out-of-sample predictive power of our model.

For crypto assets, we apply a 50 percent shock to ETH and WBTC. For stablecoins, we apply \latestrev{a 100 percent price decline to DAI and USDC as an extreme loss-of-value scenario}. The results show that standard pool indicators capture only part of the risk. In the crypto liquidity exercise, a baseline model using only the standard pool-level indicators reported by Compound explains 23.6 percent of out-of-sample variation in reserve depletion. Adding the exposure-weighted features of \latestrev{collateral providers raises} the test \(R^2\) to 0.869. For liquidity risk, the main information is therefore not simply how levered collateral providers are, but how their positions connect a pool to the rest of the protocol. Solvency is different. Standard pool indicators explain 39.3 percent of out-of-sample variation in bad debt, while the preferred parsimonious model raises the test \(R^2\) to 0.834. The gain reflects the fact that bad debt depends not only on liquidation-driven withdrawals, but also on whether affected users retain enough collateral value to cover their obligations after the shock. The stablecoin results reinforce this point. \latestrev{For DAI and USDC, adding borrower characteristics improves prediction (see Section~\ref{sec:results_solvency_stablecoin}).} The preferred model achieves a test \(R^2\) of 0.929, showing that the relevant side of the network depends on the economic role of the shocked asset.

 \rev{The \latestrev{standard pool-level} indicators the protocol publishes fail to identify the \latestrev{pool-days} most at risk. For liquidity risk, only $2.5$ percent of those they flag turn out to be among the worst hit \latestrev{in the simulated test outcomes}, against $10$ percent for a random ranking. Measures of the connectivity and concentration of a pool's counterparties raise this share to $60$ percent (Section~\ref{sec:results_incremental}).} Pool size, utilization, collateral factors, and aggregate liquidity are useful, but they largely describe pools in isolation. They do not capture how users connect collateral pools to borrowing pools, nor how those connections determine which positions are liquidated after a price shock. The relevant exposure is shaped not only by the characteristics of the pools but also by the composition and connectivity of their user base. Pools with highly leveraged or highly connected counterparties are more prone to transmitting stress across the protocol. These dynamics explain why risk monitoring in DeFi cannot rely solely on pool size, utilization, or collateral factors, but must also identify \latestrev{the users and positions connecting those pools}.



\begin{comment}
We find that the predictors of systemic risk differ 
meaningfully across asset types and risk categories, in ways that reflect the structural asymmetry between how crypto assets and stablecoins are used within the protocol. Crypto assets 
serve predominantly as collateral, so their risk profile is 
shaped by the connectivity and leverage of those who post them: 
for liquidity risk, the dominant predictors are collateral 
utilization, inflow concentration, and the size of collateral 
buffers relative to debt, all of which amplify the scale of 
liquidation-driven reserve depletion when prices fall. For 
solvency risk in crypto pools, network centrality plays a 
comparatively larger role alongside these collateral-side 
features, with PageRank emerging as the dominant own-pool 
predictor, reflecting the tendency of central pools to attract 
larger and more leveraged positions that generate unrecoverable 
shortfalls under severe shocks. \julnote{Stablecoins, by contrast, serve 
as the borrowing leg of leveraged positions rather than as 
collateral, so when a de-pegging event occurs the transmission 
channel runs through borrowers rather than collateral providers: 
the concentration, leverage, and connectivity of those who 
borrowed the stablecoin, alongside own-pool collateral 
utilization and collateral factor, are the primary 
determinants of how much bad debt accumulates when the peg 
breaks.} 
\end{comment}


\medskip


\textbf{Related Work.} There is a large and growing literature on the role of networks in economics and finance, showing that network dependencies help explain how localized shocks---whether idiosyncratic or sectoral---can propagate and amplify into macroeconomic fluctuations. Theoretically, such propagation has been formalized in models of input-output linkages and financial networks \citep{acemoglu2015systemic,gabaix2011granular}, where the topology of connections and the concentration of exposures shape systemic outcomes. Empirically, production and financial network data have been used to demonstrate how shocks transmit through supplier–buyer chains \citep{carvalho2021supply} and through firm-specific demand shifts \citep{digiovanni2014firms}. A parallel strand of research in finance models interbank markets as networks, highlighting how the structure of claims and exposures influences both contagion and stability \citep{allen2000financial,elliott2014financial}. As explained in \cite{glasserman2016contagion}, these frameworks have given rise to a variety of contagion metrics, such as DebtRank \citep{Battiston2012,Battiston2015}, which quantifies how financial distress propagates through layered interdependencies before defaults occur. Systemic risk in financial networks has also been extensively studied using complex network theory \citep{caldarelli2007scale,boss2004network}. Studies have demonstrated that while increased interconnectivity can mitigate idiosyncratic risks, it also heightens systemic fragility by creating contagion pathways if exposures are clustered among certain groups of users, as in \cite{caccioli2009eroding}.

The relevance of network models has extended beyond traditional finance, finding applications in DeFi, where lending protocols such as Compound and Aave facilitate algorithmic credit intermediation. Unlike banks, DeFi protocols operate via smart contracts, enforcing collateralized borrowing without centralized oversight \citep{bartoletti2021lendingsok}. Recent research has examined the unique arbitrage mechanisms within DeFi \citep{daian2020flash,heimbach2024nonatomic}, including miner extractable value (MEV) and the role of block proposers and arbitrageurs in shaping market outcomes \citep{milionis2023automated,azar2024information}. While traditional finance relies on established intermediation chains \citep{green2007dealer,dimaggio2017value,hollifield2017bid}, DeFi introduces a transparent but fragmented market structure with novel sources of rent extraction.

One of the mechanisms in DeFi lending networks is the role of collateralization in shaping risk transmission. Most DeFi protocols require overcollateralization, and borrower solvency is continuously assessed through a health ratio---a protocol-specific metric comparing the collateral value, adjusted by a risk factor, to the value of outstanding debt. When this ratio falls below a liquidation threshold, the borrower’s position \latestrev{becomes eligible for liquidation by other participants}. While this mechanism enforces discipline, it can also destabilize the protocol during sharp price movements by triggering cascades of liquidations \citep{sun2022liquidity}, with broader spillover effects on asset prices and interlinked positions \citep{qin2021liquidations,lehar2022systemic,gudgeon2020deficrisis}.

Although prior research in traditional finance has examined the role of collateral requirements and network structure in driving liquidation dynamics \citep{elsinger2006risk,nier2007network}, less is known about how these dynamics play out in automated settings when asset values decline rapidly. While prior work has examined systemic risk in interbank networks using eigenvalue-based stability measures \citep{caccioli2018network,glasserman2016contagion}, their applicability to DeFi lending has been limited. We help fill this gap by combining a tractable propagation algorithm with real-time health factor metrics to simulate and trace the spread of financial distress in Compound. This approach complements insights from recent models of blockchain-based financial intermediation, which highlight the roles of information asymmetry, arbitrage behavior, and endogenous feedback effects in shaping outcomes in DeFi markets \citep{cong2023inclusion,lehar2023battle}.

% Our methodology extends traditional systemic risk measures to decentralized environments, offering a more comprehensive understanding of financial contagion in DeFi lending protocols.
%We build on this financial intermediation in blockchain literature by contextualizing risk propagation within the structure of Ethereum’s proof-of-stake ecosystem \citep{abadi2018blockchain,cong2019blockchain}.

Finally, our paper builds on the spatial econometrics literature \citep{anselin1988spatial, LeSage2009}. We adopt a simulation-based approach to bypass the well-known challenges of causal inference in networks \citep{Manski, GibbonsOverman2012}. Rather than analyzing actual contagion episodes, we simulate price shocks of varying magnitudes, propagate them through the network, and regress the resulting stress measures on financial and topological features. \latestrev{Our objective is predictive rather than causal: we identify which observed features are most informative about simulated financial stress under the protocol's liquidation rules.}


\medskip

\textbf{Structure of the paper.} Section \ref{sec:Compound_Design} presents the design of Compound and explains how users' and pools' balance sheets can be constructed in order to map the protocol’s financial network. Section \ref{sec:Shock_Propagation} describes the propagation of price shocks through Compound’s network and introduces a tractable algorithm that captures this transmission mechanism. Section \ref{sec:data} presents the data, while Section \ref{sec:empirical} outlines our econometric model and the strategy used to assess its predictive performance. Section \ref{sec:results} presents the empirical results, beginning with liquidity risk in \autoref{sec:results_liquidity} and then turning to solvency risk in \autoref{sec:results_solvency}. \latestrev{Section~\ref{sec:external_validity} compares later lending protocols, with institutional details in \ref{app:later_protocols}.} Finally, Section \ref{sec:conclusion} concludes, while supplementary material and robustness checks are provided in the Appendix.

\section{The Compound Protocol}
\label{sec:Compound_Design}

\medskip

\subsection{Compound's Financial Network}
\label{sec:financial_network}

Compound V2 is a decentralized protocol that operates through smart contracts. It enables users to lend and borrow tokens in a permissionless and automated way.

On the lending side, users supply tokens to the protocol, which aggregates these deposits into \textit{liquidity pools} that contain only one type of tokens. In return for their deposits, lenders receive interest-bearing tokens called \texttt{cTokens} (such as cETH or cUSDC). These tokens represent both the user’s principal and the interest accrued over time. \latestrev{Interest accrues through the exchange rate between cTokens and the underlying asset, rather than an increase in the number of cTokens.} Because \texttt{cTokens} are themselves transferable and redeemable, they function as a tokenized claim on the underlying assets. \latestrev{Redemption requires sufficient liquidity in the pool and must not leave the user below the required collateral level.}\footnote{\latestrev{See Compound's \href{https://docs.compound.finance/v2/ctokens/}{cToken documentation}, especially ``Exchange Rate,'' ``Redeem,'' and ``Transfer.''}}

On the borrowing side, users who wish to borrow must first deposit assets into the protocol, receiving \texttt{cTokens} in return, and then use these \texttt{cTokens} as collateral to borrow other assets. Compound enforces an overcollateralization mechanism to manage risk. Each asset has an associated collateral factor, which determines how much can be borrowed relative to the value of the deposited collateral. If the value of the collateral falls below a required threshold, the position becomes eligible for liquidation, meaning that other participants can repay part of the debt in exchange for a portion of the collateral at a discount. \latestrev{For example, a user deposits ETH, receives cETH, and uses that claim as collateral to borrow USDC. The cETH remains a claim on the ETH pool: it is not converted into USDC or transferred to the USDC pool. The user's account can include several collateral assets and debts, subject to an account-wide borrowing limit. That account connects the relevant pools even though the pools do not lend directly to one another.}\par

\autoref{tab:balance_sheet} summarizes the balance sheet components of users and lending pools. Users hold \latestrev{claims on pools (inter-node assets)} in the form of \texttt{cTokens} and incur \latestrev{obligations to pools (inter-node liabilities)} through their borrowings. Lending pools, in contrast, carry a combination of tokens and loans — originated through user borrowing — on the asset side, and record user deposits as inter-node liabilities. This representation abstracts from off-protocol positions, such as users' external holdings of fiat or crypto assets not supplied to Compound, allowing us to isolate the financial linkages within the protocol. Incorporating exogenous assets and their role in risk absorption is a promising direction for future research.


\FloatBarrier 
\begin{table}[ht]
\centering
\begin{tabular}{l|c|c}
\hline
\rowcolor[HTML]{EFEFEF} 
                           & \textbf{User}     & \textbf{Pool}   \\ \hline
\rowcolor[HTML]{CFE2F3} 
\multicolumn{3}{l}{\cellcolor[HTML]{CFE2F3}\textbf{Assets}}      \\
External assets            & $\emptyset$               & Tokens          \\
Inter-node assets &
  \begin{tabular}[c]{@{}c@{}}Deposits\\ (cTokens)\end{tabular} &
  \begin{tabular}[c]{@{}c@{}}Loans
  \end{tabular} \\ \hline
\rowcolor[HTML]{E6B8AF} 
\multicolumn{3}{l}{\cellcolor[HTML]{E6B8AF}\textbf{Liabilities}} \\
External liabilities       & $\emptyset$               & $\emptyset$               \\
Inter-node liabilities     & Borrows           & Deposits        \\ \hline
\rowcolor[HTML]{FCE5CD} 
\multicolumn{1}{c|}{\cellcolor[HTML]{FCE5CD}\textbf{Equity}} &
  Deposits - Borrows &
  Tokens + Loans - Deposits \\
\hline
\end{tabular}
\vspace{3mm}
\caption{Balance sheet structure of lending pools and users.}
\label{tab:balance_sheet}
\end{table}

Using this asset-liability decomposition, we can represent Compound's financial network as a directed bipartite graph, where the nodes consist of lending pools $\mathcal{P} = {1, ..., p}$, and users $\mathcal{U} = {p+1, ..., p+u}$. An edge of weight $L_{ij}$ denotes a liability of node $i$ to node $j$---that is, node $i$ owes $L_{ij}$ to node $j$. 
Due to the bipartite nature of the network, the liability matrix $L$ satisfies $L_{ij} = 0$ for all $i, j \in \mathcal{U}$ or $i, j \in \mathcal{P}$. Liabilities only exist between users and pools; not within the same group.
%we set $A_{ij} = 0$ if both $i$ and $j$ belong to the same set ($\mathcal{P}$ or $\mathcal{U}$).

\autoref{fig:Bipartite_graph} provides an illustrative example of the bipartite financial network generated by the Compound protocol. Directed edges represent financial exposures, with arrows pointing from debtors to creditors. Each node is associated with a stylized balance sheet decomposing assets and liabilities. For clarity, node sizes are not drawn to scale. The figure highlights the dual role users play in the system: some users act as pure lenders (e.g., the second user from the left), supplying assets to lending pools without initiating any borrowings, while others both supply and borrow. This intermediation is facilitated by the pools, which aggregate deposits and extend credit to users against collateral. In the pools’ balance sheets, deposits appear as liabilities, while outstanding loans appear as inter-node assets. For users, the mirror image holds: \texttt{cTokens} deposits are assets, and borrowings are inter-node liabilities. By representing users and pools as distinct, but financially interlinked agents, this bipartite structure captures the exposure patterns that arise endogenously and serves as the foundation for the analysis of shock propagation.

\FloatBarrier 
\begin{figure}[ht]
    \includegraphics[width=\linewidth]{figures/Bipartite_Network.png}
    \caption{Bipartite network of Compound.\\
    \scriptsize{Note: Arrows point to creditors. The second user from the left is a pure lender. Node sizes are not proportional to balance sheet sizes.}}
    \label{fig:Bipartite_graph}
\end{figure}

\medskip

\subsection{Interest Rate Model and Liquidations}

The soundness of Compound hinges on its ability to uphold two core commitments: (i) providing on-demand liquidity to lenders who wish to withdraw their tokens; and (ii) ensuring the solvency of its loans by maintaining sufficient collateral backing. The former is underpinned by the protocol’s interest rate model; the latter by its liquidation mechanism, which serves as a safeguard against undercollateralization.

\medskip

\textbf{Interest Rate Model.} The interest rates at which users can borrow or lend are determined algorithmically by the Interest Rate Model (IRM hereafter) based on supply and demand within each pool. The borrowing rate for reserve \(i\), denoted \(r^B_i\), is a function of the utilization ratio \(U_i\), which measures the share of reserves in pool \(i\) that is currently borrowed. 
% Using the notation introduced above, \(U_i\) is given by
% \begin{equation}
    % U_i = \frac{\sum_{j \in \mathcal{U}} L_{ji}}\{\sum_{j \in \mathcal{U}} L_{ij}},
% \end{equation}
% where the numerator represents total borrowings from pool \(i\) by users, and the denominator represents the total supply deposited into pool \(i\). 
The protocol must ensure that the utilization ratio does not approach one, as full utilization would prevent lenders from withdrawing their funds, thereby violating the promise of on-demand liquidity.

The IRM defines, for each reserve, a mapping from the utilization ratio \(U_{i} \in [0,1]\) to the corresponding borrowing rate \(r^B_{i} \in \mathbb{R}_+\). Interest payments are distributed to lenders proportionally to their deposits, with a fraction \(\tau_{i} \in (0,1)\), known as the reserve factor, retained by the protocol. Accordingly, the lending rate \(r^L_{i}\) is related to the borrowing rate $r^B_{i}$ through the following equation
\begin{equation*}
    r^L_{i} = (1 - \tau_{i})\cdot U_{i} \cdot r^B_{i}(U_{i}).
\end{equation*}

In practice, the mapping $r^B_{i}(U_{i})$ is piecewise linear, with a kink at an optimal utilization level $U^{\text{opt}}_{i} \in (0,1)$ determined by governance.\footnote{Formally, the borrowing rate is given by 
\begin{equation*}
    r_{i}^{B} = \begin{cases} U_{i} \cdot \dfrac{r^{\text{opt}}_{i}}{U^{\text{opt}}_{i}}, & \text{if } U_{i} \leq U^{\text{opt}}_{i}, \\[10pt] r^{\text{opt}}_{i} + \left(U_{i} - U^{\text{opt}}_{i}\right) \cdot \dfrac{r^{\text{full}}_{i} - r^{\text{opt}}_{i}}{1 - U^{\text{opt}}_{i}}, & \text{if } U_{i} > U^{\text{opt}}_{i}, \end{cases}
\end{equation*}
where $r^{\text{opt}}_{i}$ is the borrowing rate at the optimal utilization ratio, and $r^{\text{full}}_{i}$ is the borrowing rate at full
utilization $U_{i}=1$.} This functional form is designed to approximate a convex, increasing relationship while remaining computationally efficient in order to minimize gas costs.\footnote{The gas cost refers to the fee incurred when executing a transaction on the Ethereum blockchain. See \cite{Bertucci2026} for a thorough analysis of the optimal design of the IRM.} The kink determines the utilization threshold at which the IRM transitions from a moderate slope to a significantly steeper one. Below this kink, borrowing rates increase gradually with utilization. Beyond this kink, interest rates rise sharply, creating a strong economic incentive for liquidity provision.

\medskip

\textbf{Liquidations.} Since Compound's users are pseudonymous, the only way to guarantee that they have an incentive to repay their loans is by ensuring that their positions remain overcollateralized at all times. The solvency of each user $i \in \mathcal{U}$ is summarized by their \textit{health ratio} $H_i$, defined as

\begin{equation}
    H_i = \frac{\sum_{j \in \mathcal{P}} c_j\cdot L_{ji}}{\sum_{j \in \mathcal{P}} L_{ij}},
    \label{eq:Health_Ratio}
\end{equation}
where $c_j \in (0,1]$ is the \textit{collateral factor} of pool $j$, a parameter that determines what fraction of deposited assets may be borrowed against. The numerator, which weights each collateral position $L_{ji}$ of user $i$ by its collateral factor $c_j$, measures her borrowing capacity. The denominator represents the user's total outstanding debt, summing her borrowed amounts $L_{ij}$ across all pools.


Because the protocol accepts volatile assets as collateral, solvency must be actively maintained through a liquidation mechanism that targets positions at risk of becoming undercollateralized. Specifically, when the user's health ratio defined in (\ref{eq:Health_Ratio}) falls below one, the position becomes eligible for liquidation: third-party liquidators may repay a portion of the outstanding debt and, in exchange, seize a corresponding amount of the borrower's collateral at a discount, hereafter referred to as the liquidation bonus $LB$. The discounted price incentivizes timely intervention and helps restore the position toward solvency, hopefully before the protocol incurs a loss. 


To bound the magnitude of any individual intervention, the protocol enforces a \textit{close factor}, which limits the fraction of \latestrev{each borrowed asset} that may be repaid within a single liquidation event (typically up to 50\%). \latestrev{This is a cap, not a requirement to repay the maximum.}\footnote{\latestrev{Compound, \href{https://docs.compound.finance/v2/comptroller/}{Comptroller documentation}, ``Close Factor.''}} Following \latestrev{a liquidation}, the borrower’s position is recomputed. If the resulting health ratio exceeds unity, the liquidation process terminates. Otherwise, subsequent liquidation rounds may be executed on the remaining collateral. This iterative procedure continues until either the position is restored to solvency or the collateral is fully depleted, in which case the protocol incurs residual bad debt.


\section{Shock Propagation} 
\label{sec:Shock_Propagation}
\medskip

\subsection{Liquidity and Solvency Risks} 
\label{sec:Risks}

Unlike traditional financial markets, which may trigger circuit breakers to suspend operations during periods of extreme stress, smart contracts remain operational at all times. Lending protocols therefore need to be resilient to large price fluctuations, ensuring that sudden market movements do not compromise their solvency or their capacity to provide liquidity. To understand why price fluctuations threaten these two properties, it is instructive to trace how they propagate through the financial network. The left-hand panel of \autoref{fig:liquidation_panels} presents a stylized setting in which a user deposits ETH as collateral to borrow stablecoins (USDCs). The horizontal red line crossing the user's balance sheet represents the liquidation threshold, i.e., the point at which her health ratio falls to one and the position becomes eligible for liquidation.

The right-hand panel of \autoref{fig:liquidation_panels} shows how this network is reconfigured after a sudden 20\% decline in the price of ETH. The mechanical impact of this price shock is to shrink the assets and liabilities of the ETH pool proportionally, as well as the assets of the borrower who posted ETH as her sole collateral. Since the user's nominal debt is unchanged, this adjustment lowers her health ratio. In our example, the drop is sufficient to push the position below the liquidation threshold. To enhance visual clarity, we illustrate the case of a full liquidation. 

%\FloatBarrier 
\begin{figure}[H]
    \begin{subfigure}[b]{0.48\textwidth}
        \centering
        \includegraphics[width=\textwidth]{figures/Balance_Sheet_PreShock.png}
        \caption{Pre-liquidation state.}
        \label{fig:pre_liquidation}
    \end{subfigure}
    \hfill
    \begin{subfigure}[b]{0.48\textwidth}
        \centering
        \includegraphics[width=\textwidth]{figures/Balance_Sheet_PostShock.png}
        \caption{Post-liquidation state.}
        \label{fig:post_liquidation}
    \end{subfigure}
    \caption{Balance sheet impact of a liquidation event triggered by a drop in the price of ETH. Panel (a) shows the
    pre-liquidation state; panel (b) shows the post-liquidation state.}
    \label{fig:liquidation_panels}
\end{figure}

As shown in panel (b), the liquidator takes two \latestrev{linked} actions. First, it repays the borrower's outstanding USDC loan, which is recorded on the liability side of the borrower's balance sheet and on the asset side of the USDC pool. This repayment converts the user's debt into reserves, thereby lowering the utilization ratio of the USDC pool. In exchange, the liquidator \latestrev{receives the borrower's cETH and redeems it for underlying ETH}. \latestrev{Receiving cETH and redeeming it are distinct actions; the simulation assumes that the seized claims are redeemed.} This withdrawal reduces the deposits of the ETH pool, increasing its utilization ratio.

The overall effect of the liquidation is thus asymmetric across pools: it alleviates liquidity pressure in the borrowing pool while tightening it in the collateral pool. This stylized example highlights that Compound is exposed to two primary categories of financial risk, each arising from distinct but interrelated mechanisms within its liquidation framework:

\begin{enumerate}
    \item \textbf{Liquidity risk}: This risk materializes when liquidators seize and convert tokens used as collateral. In doing so, they withdraw liquidity from the corresponding collateral pools. If liquidation activity becomes sufficiently large, it can deplete the available reserves. In such a scenario, depositors may face delays or, in extreme cases, an inability to redeem their \texttt{cTokens} for the underlying asset, thereby compromising the protocol’s promise of on-demand liquidity.

    \item \textbf{Solvency risk}: When the value of a borrower’s collateral is insufficient to cover the outstanding debt plus the liquidation bonus, the liquidation mechanism cannot fully recover the borrowed amount, and the resulting shortfall is recorded as \textit{bad debt}.\footnote{See \autoref{sec:Propagation_Algo} and \cite{TovanichKWP23}.} If left unmitigated, persistent bad debt can erode confidence in the protocol by creating a mismatch between users' deposits and recoverable assets.
\end{enumerate}

Compound’s design is predicated on mitigating these two risks. To this end, the protocol’s governance can adjust the two key parameters introduced above: (i) the kink in the interest rate curve and (ii) the collateral factor. By calibrating both the position of the kink and the steepness of the post-kink slope, governance can influence the likelihood that pools reach critically high utilization levels. This mechanism addresses liquidity risk by helping ensure that sufficient reserves remain available. The collateral factor, by contrast, is central to the management of both risks. A higher collateral factor improves capital efficiency for borrowers but reduces the buffer available to absorb adverse price movements. By fine-tuning this parameter across asset classes, governance can adjust the probability that liquidations deplete reserves and generate bad debt. 

These parameters form a complementary risk management framework: the IRM primarily regulates short-term liquidity dynamics, while the collateral factor governs the structural resilience of the protocol against both illiquidity and insolvency.

\medskip

\subsection{Propagation Algorithm}
\label{sec:Propagation_Algo}

Having established the central role of price shock containment in the management of protocol's risk, we now introduce a simulation algorithm that goes beyond qualitative insights to quantify both liquidity and solvency risks. To ensure tractability of the propagation mechanism, we impose two simplifying assumptions:

\begin{enumerate}[label=\textbf{A\arabic*}]
    \item \label{hyp:H1} \textbf{Liquidation sizes}: Liquidated amounts are set such that, whenever feasible, the borrower’s health factor is restored exactly to one.

    \item \label{hyp:H2} \textbf{Multi-asset positions}: For users with collateral or liabilities across multiple pools, liquidations are assumed to be executed simultaneously across all pools, with liquidation shares proportional to the relative size of liabilities in each pool.
\end{enumerate}

\latestrev{Assumption \ref{hyp:H1} can understate liquidation volumes relative to repayment up to the close-factor cap.} In Compound, liquidators may \latestrev{repay up to 50\% of an individual borrowed asset} even when a smaller amount would suffice to restore a position to a healthy level. This rule \latestrev{permits a discontinuity} at the liquidation threshold, whereas our framework generates smooth liquidation profiles.\footnote{\latestrev{Aave V4 also uses a target-health-factor approach, although its target, liquidation bonus, and residual-position rules differ from ours; see \ref{app:later_protocols_aave4}.}} \latestrev{The comparison in Section~\ref{sec:liquidation_rule} supports a conservative approximation at the benchmark shock: the close-factor outcome is 7.5\% higher. The gap narrows for larger shocks, and bad debt is identical under the two rules with the other simulation assumptions held fixed.}

Assumption \ref{hyp:H2} abstracts from the fact that liquidators may selectively target a subset of pools to minimize gas costs and manage risk exposure. However, because pool selection depends on idiosyncratic liquidator strategies, explicitly modeling such behavior would require a detailed characterization of these decision rules. We leave this promising avenue for future research and instead adopt an agnostic specification that treats all pools symmetrically.

Relaxing either assumption would require agent-based methods (e.g., \cite{TovanichKWP23}), which introduce stochastic variation into simulations and complicate the identification of robust predictive features. Our approach instead strikes a pragmatic balance between realism and tractability by abstracting from the behavioral heterogeneity of liquidators, whose explicit modeling would require a dedicated structural framework.

To identify positions eligible for liquidation, we first compute total collateral and total debt for each user $i \in \mathcal{U}$:
\begin{equation}
\label{eq:CiDi}
C_i = \sum_{j=1}^p L_{j i}, \qquad D_i = \sum_{j=1}^p L_{i j},
\end{equation}
where $p=|\mathcal{P}|$ is the number of pools. We then define \textit{pledgeable collateral} as the collateral-factor-weighted sum of collateral positions
\begin{equation}
    CL_i = \sum_{j=1}^p c_j \cdot L_{j i}.
    \label{eq:CLi}
\end{equation}
Since the health factor is given by $H_i = CL_i / D_i$, user $i$ is eligible for liquidation whenever $CL_i < D_i$.

Let $LB$ denote the liquidation bonus for liquidators. We assume throughout that
\begin{equation}
    \label{eq:solvability_condition}
    \max_{j}\{c_j\} \cdot (1+LB) < 1,
\end{equation}
since, as discussed below, it is a necessary condition for the solvency of the protocol.

The propagation rule depends on whether the user's position generates bad debt:
\begin{enumerate}
    \item \textbf{No bad debt}: The user's collateral is sufficient to cover her debt, i.e.,
    \begin{equation}
        \label{eq:no_bad_debt}
        C_i \geq (1 + LB) \cdot D_i.
    \end{equation}\noindent
    Condition~(\ref{eq:solvability_condition}) implies that no bad debt arises when the health factor is exactly equal to one.\footnote{For a user with $H_i = 1$, we have $CL_i = D_i$. Using~(\ref{eq:solvability_condition}) and the definition of $CL_i$, it follows that
    \[
    C_i \geq \frac{CL_i}{\max_j\{c_j\}} = \frac{D_i}{\max_j\{c_j\}} > (1+LB)\cdot D_i.
    \]}
    This condition is therefore necessary for protocol solvency: without it, positions with a health factor at or even above one would generate bad debt upon liquidation. In practice, the restriction is largely satisfied by Compound's parameters.

    Under the proportionality rule \ref{hyp:H2}, the liquidation ratio $\Delta_i$ rescales user $i$'s positions as
    \begin{equation*}
    \begin{cases}
        L_{ij}' = (1-\Delta_i)\cdot L_{ij},\\
        L_{ji}' = \left[1 - (1+LB)\cdot \Delta_i \cdot \frac{D_i}{C_i}\right] \cdot L_{ji},
    \end{cases}
    \end{equation*}
    where we denote post-liquidation values with a prime. Assumption \ref{hyp:H1} further requires that $\Delta_i$ restore the health factor to one, i.e., $CL_i' = D_i'$. This yields the closed-form expression
    \[
        \Delta_i = \frac{D_i - CL_i}{D_i \cdot\left[1 - (1+LB)\cdot \frac{CL_i}{C_i}\right]} \in [0,1].
    \]
    The liquidation ratio $\Delta_i$ starts at $0$ for users that are just eligible for liquidation and increases linearly to $1$ at the boundary of the bad debt region where all collateral is seized by liquidators.
    \item \textbf{Bad debt}: When $C_i < (1+LB)\cdot D_i$, the user's collateral is insufficient to cover her debt. All positions are then fully unwound:
    \begin{equation*}
    \begin{cases}
        L_{ij}' = 0, \quad \forall\, j \in [\![1,p]\!],\\
        L_{ji}' = 0, \quad \forall\, j \in [\![1,p]\!],
    \end{cases}
    \end{equation*}
    and the resulting bad debt allocated by user $i$ to pool $j$ is given by
    \[
        B_{ij} = \left[1 - \frac{C_i}{(1+LB)\cdot D_i}\right]\cdot L_{ij}.
    \]
\end{enumerate}

\medskip

These updating rules allow us to simulate the impact of price shocks by identifying which positions are liquidated and, if so, to what extent. Aggregating post-liquidation liabilities across users then yields the resulting pool balance sheets, including their updated reserves and accumulated bad debt. For brevity, we relegate the aggregation procedure to \ref{app:Aggregation} and turn to the empirical implementation of the propagation algorithm.


\section{Data}
\label{sec:data}

\medskip

\subsection{Data Description}

Our analysis relies on daily snapshots of the Compound V2 protocol collected between January 1, 2020, and June 15, 2024. The data is sourced from The Graph, which provides an indexed subgraph of the Compound V2 smart contracts. After June 2024, this hosted subgraph was deprecated following The Graph’s migration to a decentralized network.\footnote{See: \url{https://thegraph.com/blog/sunrise-here-the-graph-network/} and \url{https://thegraph.com/docs/en/archived/sunrise/}} However, by that time, Compound V2 was effectively deprecated as well, with almost negligible user activity remaining.

Each daily snapshot is taken at the final block of the day in Coordinated Universal Time (UTC) and contains comprehensive information on both lending pool characteristics and user positions:
\begin{itemize}
    \item \textbf{Lending pool characteristics:} For each lending pool, we collect snapshot data on total deposits, total borrows, available liquidity (cash), reserves, and collateral factor. The pool’s external assets are defined as the sum of cash and reserves, as shown in \autoref{tab:balance_sheet}.
    %Additionally, we record relevant policy parameters, including collateral factors and asset prices, in U.S. dollars.
    \item \textbf{Price oracles:} Compound uses the token price data from Chainlink's decentralized price oracles, which serve as the protocol’s source of truth for calculating asset valuations.\footnote{Compound v2 Price Feed:  \url{https://docs.compound.finance/v2/prices/}, and Chainlink Data Feeds: \url{https://data.chain.link/feeds}}
    We collect the token price in U.S. dollars as reported at the final block of each day.
    \item \textbf{User positions:} For each user, we record their deposit and borrow positions across supported assets, i.e., the amount each user lent and borrowed from each lending pool. These values are normalized in U.S. dollars using daily token prices to construct the liability matrix $L_{ij}$ as defined in \autoref{sec:financial_network}.
\end{itemize}

%The data enables the reconstruction of lending pool and user balance sheets, as well as the dynamic liability network. It facilitates analysis of how systemic risk propagates across different price-shock scenarios.

\medskip

\subsection{Descriptive Statistics}


We examine the evolution of Compound V2 using the aggregate lending and borrowing positions of users, as shown in \autoref{fig:lending_pools}. The protocol experienced its peak activity in 2021, followed by a gradual decline in usage over subsequent years. The figures in the left and middle panels reveal persistent user preferences: ETH and USDC consistently serve as the dominant supplied assets, while borrowing activity remains primarily concentrated in USDC and DAI. On the right panel, the total equity of the protocol increased over time, reflecting the accumulation of profits retained in the pool reserves.

\FloatBarrier 
\begin{figure}[ht]
    \includegraphics[width=\linewidth]{figures/lending_pools.pdf}
    \caption{Evolution of deposit and borrow amounts of the top lending pools in Compound V2.}
    \label{fig:lending_pools}
\end{figure}

\autoref{fig:users} documents additional trends in the number of active users, health ratios, self-borrowing prevalence,\footnote{Self-borrowing refers to the practice of depositing an asset as collateral and borrowing the same asset from the same pool. Users engage in this strategy primarily to extract rewards from protocol incentives.} and utilization ratios over time. To focus on economically significant users, we consider only lenders and borrowers with positions exceeding 100 USD. The number of active participants grew steadily from early 2020, reaching a peak of approximately 35,000 lenders and nearly 5,000 borrowers by the end of 2021. Both figures declined sharply in mid-2022 and stabilized after 2023, albeit at much lower levels of total value locked.

\FloatBarrier 
\begin{figure}[ht]
    \includegraphics[width=\linewidth]{figures/users.pdf}
\caption{Evolution of Compound V2: lenders and borrowers, health ratio, self-borrowing and utilization ratios.\\
\scriptsize{Note: \textit{Left panel}: number of active lenders (blue)
and borrowers (red) with positions exceeding \$100.
\textit{Middle panel}: median (solid) and exposure-weighted average
(dashed) health ratio of borrowers with positions exceeding \$100,
with the shaded band showing the interquartile range (25th--75th
percentile).
\textit{Right panel}: share of self-borrowing (blue) and utilization
ratio (red).}}
    \label{fig:users}
\end{figure}

The sharp increase in borrowing activity coincides with the launch of the COMP incentive program in June 2020, which distributed rewards to users engaged in borrowing~\citep{park2023phantom}. This led to a sudden surge in both self-borrowing and the utilization ratio. However, as the level of incentives declined, self-borrowing dropped to nearly zero, and the utilization ratio stabilized between 20\% and 30\% from 2022 onward.


The middle panel of the figure displays the median and weighted average of the health ratio of borrowers, along with interquartile bands. Following the introduction of the incentive program, the median health ratio hovered around 1.5, suggesting that users maintained their positions near the liquidation threshold. The weighted average health ratio was even lower than the median during 2021 and 2022, implying that large borrowers were more aggressively optimizing their leverage positions. As self-borrowing diminished after 2022 (right panel), the weighted average health ratio increased, reflecting the exit of large, highly leveraged participants from the protocol.

\medskip

\boldpara{The top 5 pools.} Lending and borrowing activity is distributed across a set of pools authorized by the protocol’s governance. The number of active pools has increased over time, ranging from 6 pools at the protocol’s inception in May 2019 to a total of 19 pools.\footnote{See Leshner, R. (2019), ``Compound v2 is Live,'' \url{https://medium.com/compound-finance/compound-v2-is-live-157db0b7cfc8}} 
However, deposits remain highly concentrated within a small subset of pools. Figure \ref{fig:Pool_Contribution} shows that the three stablecoin pools (DAI, USDC, and USDT) and the two primary crypto pools (BTC and ETH) consistently account for more than 90\% of the total value locked in Compound. Among these pools, USDT cannot be used as collateral since its collateral factor is set to zero. We therefore focus on the other four pools since they are the only ones propagating losses induced by liquidations.

\FloatBarrier 
\begin{figure}[ht]
    \includegraphics[width=\linewidth]{figures/Pool_contribution.png}
    \caption{Share of Total Value Locked in Compound for the Top 5 Pools.}
    \label{fig:Pool_Contribution}
\end{figure}

\medskip

\boldpara{Crypto vs stablecoin pools.} In addition to being concentrated within a small subset of pools, lending and borrowing activity is also segmented between crypto-assets and stablecoins. More precisely, the vast majority of loans are originated to borrow stablecoins against crypto-assets used as collateral. In other words, most borrowers enter into repurchase-like agreements in which the cash leg is replaced by a transfer of stablecoins.

The dominance of this trading pattern is illustrated in the left panel of \autoref{fig:Share_Borrow}, which shows that, at any given point in time, between 80\% and 90\% of all outstanding loans in Compound are used to borrow one of the three major stablecoins.\footnote{See \cite{TovanichKWP23} for detailed evidence on the widespread use of crypto-assets as collateral for stablecoin borrowing.} The right panel of \autoref{fig:Share_Borrow} corroborates this finding by showing that the utilization ratio—defined as the share of loaned-out tokens relative to the total reserves of each pool—is systematically higher for stablecoin pools. Particularly striking is the very low utilization ratio of ETH and BTC pools, which hovers around 10\%, highlighting the limited demand for borrowing these assets. As we will show in subsequent sections, these marked differences have important implications for the nature of the risks associated with stablecoins versus crypto-assets, thereby motivating a separate analysis for the two asset classes.

\FloatBarrier 
\begin{figure}[ht]
    \includegraphics[width=\linewidth]{figures/Pct_UtilRatio_Borrow.png}
    \caption{Share of total borrowing and utilization ratios by asset class.}
    \label{fig:Share_Borrow}
\end{figure}

\section{Empirical Analysis}
\label{sec:empirical}

\medskip

\boldpara{Overview.} The empirical strategy proceeds in three steps. First, we simulate sudden token-price shocks of size $s$ applied separately to each of the top four pools and propagate them through the liability network using the algorithm of~\autoref{sec:Propagation_Algo}. Second, for every pool $j$ and date $t$, we extract two outcomes from the simulation output: liquidity risk, $\ell_{t,j}$, defined as the share of pool assets withdrawn by liquidators, and solvency risk, $b_{t,j}$, defined as the share of pool assets left as bad debt once positions are unwound. Third, we treat $\ell_{t,j}$ and $b_{t,j}$ as separate target variables and regress each on a panel of pool-level and depositor-level characteristics through a Spatial Lag of X (SLX) model presented in \autoref{sec:Regression_Specification}. The features that enter the regressors are defined in~\autoref{sec:features}, and the LASSO-based estimation and out-of-sample evaluation procedure is described in~\autoref{sec:modelling_strategy}.


\subsection{Regression Specification}
\label{sec:Regression_Specification}
\medskip

\boldpara{Simulated outcomes.} We simulate sudden price shocks, denoted $s$, of varying magnitudes (ranging from $10\%$ to $100\%$) applied daily to each of the top five pools. For every scenario, we construct two outcome measures that capture exposure to liquidity and solvency risk, respectively.

For a given shock size, liquidity risk is defined as:
\begin{equation}
    \ell_{t,j}
    \;=\;
    \frac{T_{t,j} - T'_{t,j}}{\textit{Assets}_{t,j}},
    \label{eq:outcome}
\end{equation}
where $T_{t,j}$ denotes the value of pool $j$'s reserves immediately after the price shock but prior to any liquidations, and $T'_{t,j}$ denotes the reserve value after both the shock and the ensuing liquidations. The numerator captures the net destruction of reserves attributable solely to the liquidation process. Normalizing by the total \(\textit{Assets}_{t,j}\) of pool \(j\), defined as the sum of reserves and loans in the pool,\footnote{See \ref{app:Aggregation} for an explicit decomposition of the pools' balance sheets.} enables meaningful comparisons across pools of heterogeneous size. By construction, $\ell_{t,j}$ belongs to $[0,1]$, with larger values indicating greater relative liquidity losses. 

Our second outcome of interest is solvency risk, defined as:
\begin{equation}
    b_{t,j}
    \;=\;
    \frac{BD_{t,j}}{\textit{Assets}_{t,j}},
    \label{eq:bad_debt}
\end{equation}
where $BD_{t,j}$ denotes the total bad debt generated in the financial network following the shock to pool $j$.\footnote{See equation (\ref{eq:Def_BD}) in \ref{app:Aggregation} for a description of how $BD_{t,j}$ is computed.} This measure isolates losses that materialize when the liquidation mechanism fails to fully recover the value of undercollateralized positions. Higher values of $b_{t,j}$ therefore reflect more severe solvency stress, as a larger share of loans extended by the protocol becomes unrecoverable. Both outcome variables are stationary: Augmented Dickey-Fuller tests reject the null of a unit root for both ETH and WBTC at $s = 50\%$.\footnote{ADF statistics for $s=50\%$: liquidity risk, $-6.34$ for ETH 
($p<0.001$) and $-2.87$ for WBTC ($p=0.049$); solvency risk, 
$-3.17$ for ETH ($p=0.022$) and $-3.40$ for WBTC ($p=0.011$). 
Lag length selected by AIC.}

%Because the empirical exercise is predictive rather than causal, we do not require the outcome variables to satisfy the assumptions associated with a dynamic time-series model. The relevant validation criterion is out-of-sample predictive performance under the chronological train/test split described below. For the main crypto specification at $s=50\%$, Augmented Dickey-Fuller tests reject the null of a unit root for both liquidity and solvency risk for ETH and WBTC.\footnote{For the main crypto specification at $s=50\%$, the ADF statistics are as follows: for liquidity risk, $-6.34$ for ETH ($p<0.001$) and $-2.87$ for WBTC ($p=0.049$); for solvency risk, $-3.17$ for ETH ($p=0.022$) and $-3.40$ for WBTC ($p=0.011$). Lag length is selected by AIC.}

\medskip
\boldpara{Relationship between simulated and realized stress.} We rely on simulated outcomes because estimating a predictive model of realized liquidation events is not feasible given the rarity of large price crashes. Stable variable selection and reliable out-of-sample performance require samples of several hundred observations. In practice, however, large crypto price crashes occur at most a few times per year, making it impossible to assemble a test set of sufficient size to evaluate predictive accuracy with any statistical confidence. The simulation-based approach circumvents this limitation by generating a rich panel of stress scenarios from the observed daily network. %\julnote{We rely on simulated outcomes because estimating a predictive model of realized liquidation events is not feasible given the rarity of large price crashes. Stable variable selection and reliable out-of-sample performance require samples of several hundred observations.H However, in practice, large crypto price crashes occur at most once or twice per year, making it impossible to assemble a test set of sufficient size to evaluate predictive accuracy with any statistical confidence. The simulation-based approach circumvents this limitation by generating a rich panel of stress scenarios from the observed daily network.}

A natural question is whether the simulated outcomes are informative about real-world financial stress. Our static simulation design abstracts from two countervailing channels that operate in real time. On the one hand, price declines are often gradual rather than instantaneous, giving proactive users time to top up their collateral or reduce their borrowing, thereby avoiding liquidation. On the other hand, intra-day price volatility can trigger successive waves of liquidations within the same trading day, amplifying total losses beyond what a single end-of-day snapshot would suggest. Because these two channels operate in opposite directions, our simulations may over- or underestimate realized liquidation volumes depending on which effect dominates in a given episode.

%We assess the distance between the simulated and realized liquidation volumes in \ref{app:data_model}. We find that inter-day price shocks must be multiplied by $1.75$ to $2$ to maximize full-sample correlation, reflecting amplification from intra-day  price fluctuations omitted by the model. Since our model targets systemic stress, we also evaluate its performance during two major crashes --- the Bitcoin sell-off of May~19, 2021 and the FTX-induced crash of November~9, 2022 --- showing that episode-specific multipliers of $1.75$ and $1.50$ allow the model to closely match realized liquidation volumes, with the peak correctly identified on the day of each crash. This close fit across both episodes lends empirical support to the model's ability to provide a useful benchmark for the scale of realized losses.

We assess the quantitative distance between simulated and realized liquidation volumes in \ref{app:data_model}, using Precision-Recall analysis and episode-level comparisons across the most extreme liquidation days in our sample.  Since our 
end-of-day simulations abstract from intra-day price dynamics, 
we apply a scalar multiplier $m \geq 1$ to the observed price 
drop to account for intra-day amplification. Treating the 
simulated liquidation-to-total-assets ratio as a continuous 
risk score, the benchmark multiplier $m = 1.50$ achieves a 
PR-AUC of $0.309$, representing a $31\times$ improvement over 
a random classifier and confirming that the simulation has 
meaningful signal for identifying extreme liquidation days. Among the largest crashes occurring in the data, the model matches the COVID crash (March 2020), 
the Bitcoin flash crash (May 2021), the January 2022 market 
correction, the Terra/Luna collapse (May--June 2022), and the 
FTX crisis (November 2022), with the peak correctly identified 
on or near the day of each crash.

These results support the interpretation of the simulation as 
isolating the structural sensitivity of the network: by holding user behavior fixed and applying a clean price shock, the simulation isolates how much of the realized stress is attributable to the topology and leverage of the network itself, independently of behavioral responses. This structural sensitivity is informative for risk monitoring, where the primary concern is identifying vulnerable configurations rather than replicating the exact dynamics of a specific episode.


\medskip

\boldpara{Econometric model.} Having defined the targets, we now specify the regression model that links them to pool- and user-level characteristics. The propagation mechanism described in Section~\ref{sec:Shock_Propagation} suggests a natural structure. When a token price falls, it is the positions of depositors who have pledged that token as collateral that become undercollateralized and can trigger liquidation cascades. The scale of the resulting liquidity losses and bad debt therefore depends not only on the pool's own balance sheet, but on the financial condition and network connectivity of the users linked to it: how leveraged they are, how close to the liquidation threshold, and how broadly their positions extend across the rest of the protocol.
% A pool connected to highly leveraged, widely connected depositors faces amplified exposure even if its own balance sheet looks sound.


We capture this counterparty channel through a \emph{Spatial Lag of X (SLX) 
model}~\citep{Elhorst2014}, which augments pool-level covariates with 
exposure-weighted averages of user characteristics:
\begin{equation}
Y_t = \mathbf{X}_t \beta + \mathbf{W}^{D}_t \mathbf{Z}_t \theta^{D} + 
\mathbf{W}^{B}_t \mathbf{Z}_t \theta^{B} + \mu + u_t,
\label{eq:SLX}
\end{equation}
where $Y_t \in \mathbb{R}^{N}$ stacks the chosen risk outcome---either 
$\ell_{t,j}$ from~\eqref{eq:outcome} or $b_{t,j}$ 
from~\eqref{eq:bad_debt}---across the $N = |\mathcal{P}|$ lending pools. 
Equation~\eqref{eq:SLX} is estimated separately for each outcome. The matrix 
$\mathbf{X}_t \in \mathbb{R}^{N \times K}$ collects $K$ pool-level financial 
and topological characteristics (groups 1.1 and 1.2 in 
Section~\ref{sec:features}). $\mathbf{Z}_t \in \mathbb{R}^{M \times K}$ is 
the matrix of time-varying user characteristics (groups 2.1 and 
2.2) at time $t$, where $M = 
|\mathcal{U}|$ denotes the number of users.\footnote{Note that the number of active users varies over time. For 
notational consistency, we define $M = |\mathcal{U}|$ as the maximum number 
of users observed at any time $t$. Accordingly, the matrices $\mathbf{Z}_t 
\in \mathbb{R}^{M \times K}$, $\mathbf{W}_t^{D} \in \mathbb{R}^{N \times M}$, 
and $\mathbf{W}_t^{B} \in \mathbb{R}^{N \times M}$ are constructed with $M$ 
rows or columns, with entries corresponding to inactive users set to zero as 
needed.} The products $\mathbf{W}^D_t \mathbf{Z}_t$ and $\mathbf{W}^B_t 
\mathbf{Z}_t$ aggregate depositors' and borrowers' characteristics by 
weighting each user's contribution by their relative exposure to the pool. 
The weight matrices are row-normalized, with entries
\begin{equation}
w^{D}_{ij,t} = \frac{L_{ij,t}}{\sum_{k} L_{ik,t}} \quad \text{and} \quad 
w^{B}_{ij,t} = \frac{L_{ji,t}}{\sum_{k} L_{ki,t}},
\end{equation}
reflecting the share of pool $i$'s liabilities owed to depositor $j$ and the 
share of pool $i$'s assets lent to borrower $j$ at time $t$. By construction, 
$\mathbf{W}^B_t$ captures the relative importance of each borrower to pool 
$i$, while $\mathbf{W}^D_t$ reflects the relative importance of each 
depositor. The coefficient vector $\beta$ in~\eqref{eq:SLX} measures the 
marginal effect of the pool's characteristics on $Y_t$, while $\theta^D$ and 
$\theta^B$ capture the marginal effects of the exposure-weighted depositor and 
borrower covariates, respectively. The vector $\boldsymbol{\mu} \in 
\mathbb{R}^{N}$ collects pool fixed effects that absorb time-invariant 
heterogeneity across pools, such as differences in collateral factor design or 
governance history, and Newey--West HAC standard errors correct for residual 
serial dependence in the disturbance vector $\mathbf{u}_t \in \mathbb{R}^N$.


% The specification focuses on the depositor side and excludes borrower characteristics as a baseline choice. This is motivated directly by the propagation algorithm in Section~\ref{sec:Propagation_Algo}: a price shock impairs the collateral of depositors, not the obligations of borrowers, so it is depositor connectivity and leverage that determine whether an initial devaluation escalates into a liquidation cascade and generates treasury losses or bad debt. Robustness checks in Appendix~\ref{app: Robustness_Borrower} confirm that adding borrower characteristics does not improve predictive performance for crypto pools, validating this parsimonious choice. 

% For stablecoin pools, where the relevant shock is a de-pegging event that directly impairs outstanding loans rather than collateral, the transmission channel runs through borrowers rather than depositors; borrower characteristics are therefore included in that specification and are retained by LASSO.}

%The SLX framework is well suited to our setting for two reasons that go beyond convention. First, it avoids the reflection problem emphasised by~\cite{Manski}: because propagation losses are computed independently for each pool via simulation, there is no simultaneity in $Y_t$ across pools, so the identifying assumption of the SLX model holds by construction.\footnote{While our SLX specification avoids simultaneity from endogenous spatial lags of the dependent variable (i.e., no $\mathbf{W}Y_t$ terms), endogeneity of covariates or their spatial lags ($\mathbf{W}X_t$) may still be a concern. Following~\citet{HalleckVegaElhorst2015}, we treat propagated losses as exogenously computed. Their ex ante simulation, independent of realized outcomes, mitigates concerns about reverse causality.} Second, unlike spatial autoregressive or Spatial Durbin models, SLX does not require inverting a spatial multiplier matrix---an operation that would be both computationally demanding and difficult to interpret in a network with daily-varying topology and thousands of active users~\citep{Elhorst2014}. In this sense, the liability network plays the role that geographic distance plays in traditional spatial models: it defines who influences whom, with exposure-weighted financial links replacing physical proximity as the channel of interaction~\citep{BenAbdallah2022,paelinck1979spatial}.

The SLX framework is well suited to our setting for two 
reasons. First, it avoids the reflection problem emphasized 
by~\cite{Manski}: because propagation losses are computed 
independently for each pool via simulation, there is no 
simultaneity in $Y_t$ across pools, so the identifying 
assumption of the SLX model holds by 
construction.\footnote{Because the liability network is jointly shaped by user behavior, protocol incentives, and market conditions, the exposure-weighted regressors ($\mathbf{W}Z_t$) are interpreted descriptively rather than causally. Our objective is prediction rather than identification of structural peer effects.} Second, SLX 
produces directly interpretable coefficients: $\theta$ 
measures the marginal effect of exposure-weighted depositor 
characteristics on pool-level risk without requiring the 
inversion of a spatial multiplier matrix, an operation that 
would be both computationally demanding and difficult to 
interpret in a network with daily-varying topology and 
thousands of active users~\citep{Elhorst2014}. Hence, the liability network plays the role that geographic 
distance occupies in traditional spatial models: it defines 
who influences whom, with exposure-weighted financial links 
replacing physical proximity as the channel of 
interaction~\citep{BenAbdallah2022,paelinck1979spatial}.

\subsection{Features Extraction}\label{sec:features}

We now describe the variables that populate $\mathbf{X}_t$ and $\mathbf{Z}_t$ in equation~\eqref{eq:SLX}. A token-price shock affects systemic risk through two linked channels. The first is the pool-level balance sheet channel: pools differ in utilization, equity buffers, collateral parameters, and the concentration of claims. The second is the counterparty channel: the users connected to a pool differ in leverage, health ratios, collateral usage, and network connectivity. These user characteristics determine whether a shock triggers liquidations, how large those liquidations are, and how losses are transmitted across pools. We organize the predictors into four groups, which map directly into the regression: groups 1.1 and 1.2 enter as $\mathbf{X}_t$, while groups 2.1 and 2.2 enter as $\mathbf{Z}_t$ before being weighted by $\mathbf{W}^D_t$ and $\mathbf{W}^B_t$.

\begin{enumerate}
\item \textbf{Pool characteristics} (populate $\mathbf{X}_t$):
\begin{enumerate}[label=\arabic{enumi}.\arabic*.]
    \item \textbf{Financial features} summarize the pool's balance sheet and protocol-defined risk parameters. These include the collateral factor, utilization ratio, equity share, pool size, the share of supplied assets used as collateral, and the self-borrowing rate. These variables capture the pool's capacity to absorb withdrawals, liquidations, and bad debt following a price shock.
    \item \textbf{Topological features} describe the pool's position in the liability network. We include measures of connectivity, such as in-degree and out-degree, concentration measures based on Herfindahl-Hirschman indices of incoming and outgoing exposures, and centrality measures such as PageRank and hub scores. These variables capture whether a pool is broadly connected or dependent on a concentrated set of users.
\end{enumerate}
\item \textbf{User characteristics} (populate $\mathbf{Z}_t$, then aggregated to the pool level via $\mathbf{W}^D_t$ and $\mathbf{W}^B_t$):
\begin{enumerate}[label=\arabic{enumi}.\arabic*.]
    \item \textbf{Financial features} summarize the condition and behavior of users connected to each pool. These variables are constructed as exposure-weighted averages at the pool level and include users' health ratios, leverage, collateral usage, and related balance sheet measures. They capture how close connected users are to liquidation and how strongly their positions depend on the shocked asset.
    \item \textbf{Topological features} describe the connectivity and exposure concentration of users connected to each pool. These variables are also aggregated to the pool level using exposure weights and include user in-degree, out-degree, flow weights, and the corresponding Herfindahl-Hirschman concentration indices. They capture whether a pool is connected to users whose positions link many parts of the protocol or to users whose exposures are more isolated.
\end{enumerate}
\end{enumerate}

In short, these four groups constitute the initial set of $P = 40$ candidate variables submitted to the estimation procedure. Full variable definitions are provided in~\autoref{tab:node_features} in \ref{app:feat_des_app}.


% ════════════════════════════════════════════════════════════════
% 4.5  Predictive Modelling Strategy
% ════════════════════════════════════════════════════════════════
\subsection{Predictive Modelling Strategy}
\label{sec:modelling_strategy}

%cm edits

Each daily network snapshot provides a cross-sectional observation of Compound’s liability network at the end of a trading day. Lending pools are observed repeatedly over time, so the data have a panel structure. Our empirical objective is not to estimate dynamic causal effects. We use the repeated daily snapshots to examine how contemporaneous pools' and exposure-weighted users' characteristics predict simulated liquidity and solvency risk. The estimates should therefore be interpreted as predictive associations that identify which features are most informative about systemic vulnerability under the propagation algorithm.

To evaluate predictive performance, we use a strict chronological train/test split, LASSO variable selection, and expanding-window (rolling-origin) cross-validation for penalty tuning. This design reflects the forecasting problem faced by risk managers: whether information available in the current network can predict simulated liquidity losses and bad debt under large token-price declines. The chronological split ensures that models are trained on past observations and evaluated on future observations, while LASSO 
penalizes model complexity and reduces overfitting. Coefficient estimates should be read as measures of predictive association, not as causal effects. \ref{app:rf_benchmark} confirms that the linear 
specification is not restrictive: the preferred OLS 
model achieves comparable accuracy than random forest benchmarks, while 
preserving the interpretability of the estimated 
coefficients.
\medskip

\boldpara{Sample splitting.}
We construct a panel dataset integrating the features described above for the major crypto assets listed on the Compound protocol, observed daily from 1 January 2021 to 1 October 2023.\footnote{Pre-2021 observations are excluded because the protocol 
operated in a structurally different regime: the collateral factor 
averaged 0.49 in 2020 versus 0.70--0.76 from 2021 onward following 
a governance change, the borrower count was approximately five times 
smaller, and pool sizes were an order of magnitude lower. Including 
pre-2021 observations materially reduces out-of-sample $R^2$, 
confirming the structural break.} The primary outcome variables are the liquidity risk $\ell_{t,j}$ and the solvency risk $b_{t,j}$, constructed under an exogenous price decline applied separately to each asset.

Although the simulated outcomes are conditionally independent given the features, the features themselves may exhibit autocorrelation. A misspecified model will fail to fully absorb this structure, inducing autocorrelation in the residuals and violating the independence assumption underlying standard cross-validation. We therefore adopt time series validation techniques, which respect temporal ordering and provide a more reliable evaluation in this setting \citep{BERGMEIR2012192}. The sample is divided chronologically in an 80/20 ratio: the first 80\,\% of observations constitute the training set ($N_{\text{train}} = 1{,}607$ observations, covering 1 January 2021 to 15 March 2023) and the remaining 20\,\% the held-out test set ($N_{\text{test}} = 401$ \rev{pool-day} observations, covering 15 March 2023 to 1 October 2023).\footnote{The 
test period is not a low-variance regime: both mean and 
variance of $\ell_{t,j}$ and $b_{t,j}$ are substantially 
higher in the test period than in the training period 
(Levene $p < 0.001$), confirming that the high test $R^2$ 
reflects genuine predictive accuracy.} This ensures that the model is always estimated on past data and evaluated on future observations, as in a genuine forecasting exercise.\footnote{All predictors are standardized to zero mean and unit variance using moments computed on the training set only, so that coefficient magnitudes are directly comparable across predictors and no information from the test period contaminates the scaling step.} 

We adopt an expanding-window (rolling-origin) design rather than a fixed-width rolling window for two reasons. First, a rolling window would reduce the effective training sample and potentially destabilise LASSO variable selection, which already operates on a relatively limited sample. Second, since the simulated outcomes are 
conditionally independent across time, horizon-specific forecast accuracy is not a concern, removing one of the principal arguments in favour of a fixed-width rolling window \citep{TASHMAN2000437}. The remaining limitation of expanding-window evaluation is sensitivity to occurrences unique to a particular origin. We address this concern in \ref{app:Robustness_Subsampling} through a subsampling robustness check, randomly dropping 10\,\% of training observations across 50 draws. Both the selected variables and the out-of-sample $R^2$ remain largely unchanged, providing direct evidence against parameter instability over the training period.

\medskip


\boldpara{Variable selection and estimation.}
We apply LASSO to the full set of $P = 40$ standardized predictors,\footnote{Zero-variance predictors are dropped prior to estimation; $P = 40$ 
holds across all shock sizes in our sample. For shock sizes where no predictor 
survives at the significancy level $\alpha = 0.001$, the threshold is relaxed to $p < 0.05$; this 
occurs only at $s \in \{0.10, 0.20\}$ where the signal is weakest.} to identify a sparse and robust subset that generalises well to the test period. The penalty parameter $\lambda$ is selected by expanding-window (rolling-origin) time-series cross-validation with $K = 5$ folds, which respects the temporal ordering of the data by always training on past observations and validating on the subsequent block. A grid of 200 values of $\lambda$ is evaluated, and the value minimising mean cross-validation MSE across folds ($\lambda_{\min}$) is 
retained for variable selection. LASSO is used exclusively as a selection device; OLS is then refit on the selected variables with Newey--West HAC standard errors ($L = 4$ lags), removing the shrinkage bias inherent in LASSO coefficients. A final trimming step retains only predictors whose post-selection OLS coefficient is individually significant at
$\alpha = 0.001$, ensuring that the final model contains only variables with robust statistical support.

% Applied to the $s = 0.50$ shock, LASSO selects 20 out of 30 predictors; after trimming, 14 variables are retained in the liquidity risk model and 10 in the solvency risk model (rising to 15
% and 11 respectively once asset fixed effects are considered).

Performance is evaluated on the held-out test set using three metrics: the coefficient of determination $R^2$, the root mean squared error RMSE, and the mean absolute error MAE. A substantial gap between training and test $R^2$ is interpreted as evidence of overfitting.

\medskip

\boldpara{Shock size.}
The analysis is conducted separately for each shock size
$s \in \{10\%, 20\%, \ldots, 100\%\}$. \rev{The informative range for
a stress test is bounded on both sides. A decline
smaller than the overcollateralization buffer ($s \leq 20\%$)
generates little propagation because it leaves most borrowers above
their liquidation threshold. A decline large
enough to exhaust all collateral has the opposite effect, since positions
are unwound in full whatever their connections and outcomes converge
across pools. Between these two extremes, liquidations
lead to contagion that is amplified by the network structure. For crypto assets (ETH and WBTC), we therefore focus
our main results on $s = 50\%$, a severe but plausible scenario for
volatile collateral: it is consistent with the drawdowns of the Terra/Luna collapse of May
2022 and the FTX crisis of November 2022, and \ref{app:data_model}
shows that intra-day declines during major crashes, scaled by a factor
of $1.50$--$1.75$ to account for intra-day amplification, are of the
same order.}

\rev{\autoref{fig:r2_by_shock_size} confirms this reading.} It reports
test and training $R^2$ of the trimmed model for each shock size. Test
$R^2$ is low for $s \leq 20\%$, rises steeply from $s = 30\%$ onward,
peaks at $s = 50\%$, and declines thereafter even as training $R^2$
continues to rise. The decline reflects the convergence described
above: the standard deviation of $\ell_{t,j}$ peaks at $s = 50\%$
($\sigma = 0.047$) and falls to $\sigma = 0.013$ at $s = 90\%$, so the
cross-sectional heterogeneity the model is designed to capture
disappears and the model fits residual noise instead. \ref{app:shock_sizes} reports results for every shock
size. \rev{Predictive accuracy is stable across the severe-shock
range, with test $R^2$ of $0.80$ at $s = 40\%$, $0.87$ at
$s = 50\%$ and $0.83$ at $s = 60\%$.}


\FloatBarrier 
\begin{figure}[ht]
\captionsetup{justification=justified,
              singlelinecheck=false}
\includegraphics[width=\linewidth]
    {figures/shock_size.png}
\caption{Predictive performance by shock size --- trimmed model
(no asset fixed effect), liquidity risk $\ell_{t,j}$ (ETH and
WBTC).\\
\scriptsize{Note: Coral dots show test $R^2$ and blue squares show training
$R^2$ on the held-out test set. The shaded red region marks the
low-signal regime ($s \leq 20\%$) where few positions are
liquidated and the outcome has limited cross-sectional variation.
}}
\label{fig:r2_by_shock_size}
\end{figure}

For stablecoin pools (DAI and USDC), the analysis focuses on \latestrev{an extreme loss-of-value scenario} ($s = 100\%$) rather than on shocks of varying magnitudes. \latestrev{This experiment assumes that the shocked stablecoin loses all of its value; it does not represent every departure from dollar parity. We retain this scenario to examine severe stress in the stablecoin analysis} and report results for both DAI and USDC.


% ════════════════════════════════════════════════════════════════
% 5.  Results
% ════════════════════════════════════════════════════════════════
\section{Results}
\label{sec:results}

We present results for liquidity and solvency risk, reporting the variables selected by LASSO and retained after trimming, the predictive performance of the preferred model on the test set, and
an interpretation of the key coefficient estimates in terms of the economic channels identified in~\autoref{sec:Risks}.

% ────────────────────────────────────────────────────────────────
\subsection{Liquidity Risk}
\label{sec:results_liquidity}

We begin by analyzing the factors that drive liquidity risk arising from a decline in the price of one of the main tokens. As discussed in~\autoref{sec:Shock_Propagation}, this risk materializes when liquidators seize and convert tokens used as collateral. To preserve
lenders' ability to redeem their deposits, the system must prevent liquidators from fully depleting the available reserves of a given pool.

\medskip

\boldpara{Cryptos vs stablecoins.}
Token withdrawals triggered by liquidations are substantially larger in crypto pools than in stablecoin pools, reflecting the much higher share of crypto deposits pledged as collateral. \autoref{fig:reserve_changes} illustrates this dichotomy by reporting changes in pool available
reserves following a 50\,\% price drop. For each pool, we report two daily statistics: (i) the change when the price drop affects the pool's own token, and (ii) the average change when the price drop affects other top-5 tokens. 

A shock to a crypto pool's own token generates substantial reserve losses, typically around 20\% of total assets, as a large fraction of deposits serves as collateral. When other tokens are shocked instead, the impact on crypto pool reserves is comparatively modest, as the two opposing channels---inflows from repaid loans and outflows from liquidated cross-collateral positions---largely cancel out. For stablecoin pools, the dynamics are markedly different: shocks to other tokens replenish stablecoin reserves as loans are repaid, while own-token shocks sometimes have negligible impact, most notably for DAI, because stablecoins are less often pledged as collateral. Two main lessons follow: liquidity risk materializes primarily through own-token shocks, and crypto pools face greater exposure to liquidation-driven reserve depletion. Since the price shocks of primary interest for stablecoins are full de-pegging events that deplete available reserves regardless of 
propagation, we restrict the liquidity risk analysis to the two main crypto assets, WBTC and ETH.

\FloatBarrier 
\begin{figure}[ht]
\captionsetup{justification=justified,
              singlelinecheck=false}
\includegraphics[width=\linewidth]
    {figures/liquidity_risk_mean_new1.png}
\caption{Changes in pool available reserves following a price
drop of 50\,\% (ETH, WBTC, DAI, and USDC).\\
\scriptsize{Note: For each pool, the
blue solid line shows the daily change in available reserves relative to total assets
when the price drop affects the pool's own token; the red
dashed line shows the average change when the price drop
affects one of the other top-5 tokens. \textit{Upper panels} correspond to crypto pools, whereas \textit{lower panels} correspond to stablecoin pools.}}
\label{fig:reserve_changes}
\end{figure}


\medskip

\boldpara{Variable selection and predictive performance.}
We first assess whether borrowers' characteristics improve the predictive 
power of the model. The robustness checks relegated to ~\ref{app: Robustness_Borrower} show that adding them does not improve 
predictive performance. We therefore restrict the set of user covariates to 
the depositor side and exclude borrower characteristics ($P=30$). This finding is 
consistent with the logic of the propagation algorithm presented in 
Section~\ref{sec:Propagation_Algo}: price shocks impair the collateral of 
depositors, so it is their connectivity 
and leverage that primarily determine whether an initial devaluation escalates 
into a liquidation cascade that generates liquidity losses. We then 
apply LASSO, which selects 20 variables at $\lambda_{\min}$. After trimming at $\alpha = 0.001$, 14 predictors remain. The asset fixed effect for WBTC is excluded
from the preferred specification, as it is not statistically significant once pool and collateral characteristics are controlled for ($\hat{\gamma} = +0.0039$, $p = 0.337$), indicating that the
two assets do not differ systematically in liquidity losses beyond what is explained by the observable predictors.

The preferred model achieves a training $R^2 = 0.835$ and a test $R^2 = 0.868$, with $\text{RMSE} = 0.0205$ and $\text{MAE} = 0.0164$ on the held-out test set. The test $R^2$ exceeding the training $R^2$ reflects the lower variance of $\ell_{t,j}$ in the test period and confirms that the model generalises well without overfitting. A comparison of all specifications is relegated to \ref{app:model_comparison_liq}. The full OLS benchmark on all 30 predictors achieves a test $R^2 = 0.784$, substantially below the trimmed model despite using twice as many variables, confirming that LASSO-based selection improves out-of-sample performance by reducing overfitting.

\autoref{fig:liquidity_predictions} reports the HAC coefficient
estimates for the preferred model, organized by predictor group,
along with the four diagnostic panels.\footnote{See \autoref{tab:coefs_comparison} in \ref{app:model_comparison_liq} for a detailed summary of the coefficients and their predictive power.} Three channels stand out. 

\FloatBarrier 
\begin{figure}[ht]
\captionsetup{justification=justified,
              singlelinecheck=false}
\includegraphics[width=\linewidth]
    {figures/liquidity_risk_prediction_50.png}
\caption{Diagnostic plots for the preferred model ---
liquidity risk $\ell_{t,j}$ (ETH and WBTC, $s = 50\%$,
14 variables, Train $R^2 = 0.835$, Test $R^2 = 0.868$).\\
\scriptsize{Note: \textit{Upper-left}: actual against fitted values on the
training set (blue). \textit{Upper-right}: actual against
fitted values on the held-out test set (coral). \textit{Lower-left}:
HAC-estimated coefficients with 95\,\% confidence intervals,
ordered by magnitude.
\textit{Lower-right}: residuals over time for the training
(blue) and test (coral) periods, separated by the dashed
vertical line at the train/test boundary (March 2023).}
}
\label{fig:liquidity_predictions}
\end{figure}
\FloatBarrier 
The most important channel operates through the
characteristics of \emph{collateral providers}. The value-weighted
collateral utilization rate\footnote{This feature corresponds to the share of deposits used as collateral across lenders, averaged with weights proportional to the size of their deposits. See \ref{app:feat_des_app} for more details on its construction and that of other features.} carries the largest positive coefficient
($+0.0303$), indicating that pools whose collateral base is more heavily pledged sustain larger liquidity losses. The value-weighted in-flow weight concentration also carries a 
large positive coefficient ($+0.0257$), indicating that pools 
whose collateral base is more concentrated face amplified 
liquidity losses. The value-weighted health ratio carries a positive coefficient ($+0.0153$): \latestrev{higher collateral-provider health ratios are associated with greater simulated reserve depletion, conditional on the other predictors}.

The second channel operates through \emph{own-pool financial
characteristics}. Collateral utilization ($+0.0131$) and the
utilization ratio ($+0.0057$) both carry positive coefficients,
confirming the finding that pools with more actively deployed
reserves are more exposed. By contrast, the collateral factor
($-0.0122$) \latestrev{is negatively associated with simulated reserve depletion}.

The third channel operates through the \emph{cross-pool network
structure}. Greater heterogeneity in pool connectivity, measured
by in-degree dispersion, is associated with larger losses
($+0.0106$), while the value-weighted average out-degree across
pools carries a negative coefficient ($-0.0083$).

The diagnostic panels in \autoref{fig:liquidity_predictions}
confirm the model's stability. Residuals in the training period
show transient spikes during the Terra/Luna collapse (May 2022)
and the FTX crisis (November 2022), both of which were associated
with unusually large liquidation waves. In the test period,
residuals remain centred around zero with substantially lower
volatility, and the transition at the train/test boundary is
smooth, confirming structural stability across the two sub-periods.

\medskip

\boldpara{Forecast error dynamics.}
\autoref{fig:forecast_errors} examines whether prediction
accuracy improves as the model accumulates more data, using
two complementary approaches. The left and centre panels plot
absolute residuals over time from the preferred model,
for the training and test periods respectively. 
\FloatBarrier 
\begin{figure}[ht]
\captionsetup{justification=justified,
              singlelinecheck=false}
\includegraphics[width=\linewidth]
    {figures/errors.png}
\caption{Forecast error dynamics --- liquidity risk
$\ell_{t,j}$ (ETH and WBTC, $s = 50\%$).\\
\scriptsize{Note: \textit{Left}:
absolute residuals from the fixed preferred model over the
training period (blue dots), with a 30-day rolling mean and
an OLS trend line. \textit{Centre}: absolute residuals over
the test period (coral dots), with a 15-day rolling mean
and an OLS trend line. \textit{Right}: monthly mean absolute
error from a recursive expanding-window exercise in which
the full LASSO$\,\to\,$OLS pipeline is refit at the start
of each month using all data available up to that point;
the dashed line is the OLS trend.}}
\label{fig:forecast_errors}
\end{figure}
\FloatBarrier 
The right panel
reports monthly mean absolute errors from a recursive
expanding-window exercise in which the full LASSO$\,\to\,$OLS
pipeline is refit at the start of each month using all
available data up to that point.
Three findings emerge. First, training residuals show a 
statistically significant downward trend (slope $= -3 \times 
10^{-6}$, $p < 0.001$), indicating that the fixed model's 
in-sample fit improves gradually over the training period. 
Second, this improvement does not translate into better 
out-of-sample performance: the recursive expanding-window 
exercise in the right panel shows that monthly MAE does not 
decline significantly as the training set grows (slope $= -1.38 
\times 10^{-4}$, $p = 0.457$), indicating that the 
data-generating process is stable and early observations remain 
informative. The in-sample improvement therefore reflects the 
model fitting the training data more tightly over time rather 
than genuinely learning a changing relationship. Third, the test 
period residuals show an upward trend (slope $= 5.0 \times 
10^{-5}$, $p < 0.001$), but this reflects a short 7-month 
window and elevated errors concentrated in a specific episode 
toward the end of the sample rather than a systematic 
deterioration of predictive accuracy. Taken together, these 
results support the expanding-window design: since the recursive 
exercise shows no meaningful improvement from accumulating more 
data, a fixed-width rolling window that discards early 
observations would offer no advantage.


% 


\subsection{Solvency Risk}
\label{sec:results_solvency}

We now turn to the factors that drive solvency risk. As discussed in~\autoref{sec:Shock_Propagation}, this risk arises when a decline in collateral values is so severe that liquidators are unable to fully recover borrowers’ nominal debt, thereby leaving residual bad debt in the system.

\medskip

\subsubsection{Crypto Assets}
\label{sec:results_solvency_crypto}


\boldpara{Variable selection and predictive performance.} LASSO at $\lambda_{\min}$ selects 18 variables. After trimming at $\alpha = 0.001$ and including the asset fixed effect for WBTC, the preferred specification retains 10 predictors. Unlike the liquidity risk model, the WBTC fixed effect is retained here: WBTC pools generate systematically higher bad debt than ETH pools beyond what is explained by the structural covariates ($\hat{\gamma} = +0.0086$, $p < 0.001$).

The preferred model achieves a training $R^2 = 0.838$ and a test $R^2 = 0.834$, with $\text{RMSE} = 0.0072$ and $\text{MAE} = 0.0058$. The near-equality of training and test $R^2$ confirms that the
determinants of bad debt are stable across the two sub-periods.
A comparison of all specifications is reported in~\autoref{tab:model_comparison_solv}, which we relegate to~\ref{app:model_comparison_solv}. The full OLS benchmark, using all 30 predictors, achieves a test $R^2 = 0.835$, essentially identical to that of the trimmed model despite employing three times as many variables. This confirms that LASSO successfully compresses the predictive information in the full set of candidates into a parsimonious specification without sacrificing out-of-sample accuracy.

\medskip

\boldpara{Coefficient estimates.} 
\autoref{fig:solvency_predictions} reports the coefficient
estimates, organized by predictor group.\footnote{See~\autoref{tab:coefs_comparison_solv} in~\ref{app:model_comparison_solv} for a complete description of the coefficient estimates and their predictive power.} The results confirm three
channels, broadly mirroring those identified for liquidity risk but
with distinct leading predictors.

\FloatBarrier
\begin{figure}[ht]
\captionsetup{justification=justified,
              singlelinecheck=false}
\includegraphics[width=\linewidth]
    {figures/solvency_risk_predictions.png}
\caption{Diagnostic plots for the preferred model ---
solvency risk $b_{t,j}$ (ETH and WBTC, $s = 50\%$,
10 variables, Train $R^2 = 0.838$, Test $R^2 = 0.834$).\\
\scriptsize{Note: \textit{Upper-left}: actual against fitted values on the
training set (blue). \textit{Upper-right}: actual against
fitted values on the held-out test set (coral). \textit{Lower-left}:
HAC-estimated coefficients with 95\,\% confidence intervals,
ordered by magnitude. The WBTC fixed effect is the largest
single predictor.
\textit{Lower-right}: residuals over time for the training
(blue) and test (coral) periods, separated by the dashed
vertical line at the train/test boundary (March 2023).}}
\label{fig:solvency_predictions}
\end{figure}

The most important channel operates through the characteristics of \textit{collateral providers}. The value-weighted collateral utilization rate carries the largest coefficient in absolute value ($+0.0071$) as pools with more heavily pledged collateral bases generate more bad debt.
The value-weighted in-flow weight concentration ($+0.0011$) reinforces
this finding: concentrated collateral structures amplify solvency
exposure because the failure of a single large position generates
more unrecoverable shortfalls.

The second channel operates through \emph{own-pool network
position}. PageRank centrality ($+0.0069$) is the most important
own-pool predictor. Its coefficient nearly
doubles from $+0.0033$ (no FE) to $+0.0069$ (with FE), \latestrev{after accounting for the average difference between ETH and WBTC}.

The WBTC fixed effect ($\hat{\gamma} = +0.0086$, $p < 0.001$)
absorbs a residual asset-level difference in solvency exposure
not captured by the structural covariates. \latestrev{The fixed effect does not distinguish among possible explanations for that difference.}
Coefficient signs are stable across specifications with and
without the fixed effect, confirming that its inclusion sharpens
rather than distorts the estimated predictive associations
(see~\autoref{tab:coefs_comparison_solv} in~\ref{app:model_comparison_solv}).
\medskip

% ────────────────────────────────────────────────────────────────
\subsubsection{Stablecoin Pools}
\label{sec:results_solvency_stablecoin}

% \boldpara{Transmission mechanism.}
% The solvency risk analysis for stablecoin pools differs from the
% crypto case in both the nature of the shock and the transmission
% channel. In crypto pools, a price shock erodes the value of
% crypto collateral posted by borrowers, triggering liquidations
% that generate bad debt when seized collateral proves
% insufficient to cover outstanding loans. The transmission
% channel therefore runs from collateral \emph{providers} to the
% pool: it is the depositors who post collateral whose positions
% become impaired.

% In stablecoin pools, the mechanism is different. First, unlike 
% crypto assets, where variation in shock size traces out a 
% meaningful risk frontier, stablecoins either maintain their peg 
% or they do not, making partial shocks economically implausible. 
% This motivates our focus on full de-pegging scenarios ($s = 100\%$). 
% Second, because stablecoins are predominantly \emph{borrowed} 
% rather than pledged as collateral, a de-pegging event directly 
% impairs the value of outstanding loans in the pool. The 
% transmission channel therefore operates through the pool's own 
% financial structure and the network characteristics of collateral 
% providers, as in the crypto case.

%\boldpara{Variable selection and predictive performance.} In contrast to crypto pools, borrower characteristics add substantial predictive power when the model is applied to stablecoin pools. This result is intuitive: stablecoins are primarily borrowed rather than used as collateral, so the characteristics of borrowers carry fundamental information about how stablecoin pools are embedded in the financial network. 
% For consistency with the crypto analysis, the candidate predictor 
% set is restricted to pool-level and collateral-side features 
% ($P = 30$). 
%\julnote{UPDATE: LASSO at $\lambda_{\min}$ followed by trimming at $\alpha = 0.001$ retains 8 predictors. The preferred specification includes a DAI asset fixed effect, which enters negatively and significantly ($\hat{\gamma} < 0$, $p < 0.001$), indicating that DAI pools accumulate systematically less bad debt than USDC pools conditional on the covariates.}

%The preferred model achieves a training $R^2 = 0.924$ and a test $R^2 = 0.781$, with $\text{RMSE} = 0.0237$, on the held-out test set. The gap between training and test $R^2$ is larger than in the crypto case, reflecting elevated residuals concentrated in the test period (visible in \autoref{fig:solvency_stablecoins_size100}, lower-right panel), likely driven by idiosyncratic de-pegging dynamics in early 2023 that are not fully captured by the collateral-side features alone.

%\medskip

%\boldpara{Coefficient estimates.}
%\autoref{fig:solvency_stablecoins_size100} displays the diagnostic plots for the preferred model. The 8 retained predictors span own-pool and collateral network features, with two channels of particular interest.

%\FloatBarrier
%\begin{figure}[ht]
%\captionsetup{justification=justified,
%              singlelinecheck=false}
%\includegraphics[width=\linewidth]
%    {figures/summary_trim_fe_size100_stablecoins.png}
%\caption{Diagnostic plots for the selected model --- solvency risk
%$b_{t,j}$ (DAI and USDC, $s = 100\%$, 8 variables,
%Train $R^2 = 0.924$, Test $R^2 = 0.781$).\\
%\scriptsize{Note: 
%\textit{Upper-left}: actual against fitted values on the training
%set (blue).
%\textit{Upper-right}: actual against fitted values on the
%held-out test set (coral).
%\textit{Lower-left}: HAC-estimated coefficients with
%95\,\% confidence intervals, ordered by magnitude.
%\textit{Lower-right}: residuals over time for the training period 
%(blue) and the test period (coral), separated by the dashed
%vertical line at the train/test boundary (March 2023).
%}}
%\label{fig:solvency_stablecoins_size100}
%\end{figure}

%The most prominent channel operates through \emph{own-pool financial characteristics}. Collateral utilization ($+$) and the collateral factor ($+$) carry the largest positive coefficients, indicating that pools accepting more collateral at higher loan-to-value ratios accumulate more bad debt when the stablecoin de-pegs. The cross-pool self-borrowing rate ($+$) also enters positively, consistent with internally recycled positions being more exposed to a sudden loss of parity. By contrast, the value-weighted self-borrowing rate of collateral providers ($-$) enters with the largest negative coefficient in the model, suggesting that collateral providers with higher self-borrowing exposure are associated with lower bad debt accumulation, possibly reflecting their capacity to self-liquidate when the peg breaks.

%The second channel operates through \emph{collateral network features}. The value-weighted out-flow concentration and out-weight HHI of collateral providers enter negatively, while the value-weighted in-degree of collateral providers (cross-pool) also enters negatively. Together these indicate that more dispersed and less concentrated collateral structures are associated with lower bad debt under a de-pegging event, mirroring the finding from the crypto solvency analysis.

\boldpara{Variable selection and predictive performance.} LASSO at $\lambda_{\min}$ selects variables from the full candidate set of $P = 40$ predictors, which for stablecoin pools includes both collateral provider and borrower characteristics.
After trimming at $\alpha = 0.001$, 11 predictors are retained. No asset fixed effect is included, as the dummy constructed for DAI versus USDC carries a zero coefficient and an undefined $p$-value, confirming that the two stablecoin pools do not differ
systematically in bad debt exposure beyond what is captured by the covariates.

The preferred model achieves a training $R^2 = 0.958$ and a test $R^2 = 0.929$, with $\text{RMSE} = 0.0237$, on the held-out test set. Both the high level and the near-equality of training and test $R^2$ confirm that the model generalizes well to the
out-of-sample period without overfitting, and that the determinants of stablecoin bad debt under a de-pegging scenario are highly predictable from the observable network and financial characteristics of the pool and its users.

\medskip

\boldpara{Coefficient estimates.}
\autoref{fig:solvency_stablecoins_size100} displays the
diagnostic plots for the preferred model. The 11 retained
predictors span all four feature groups, with three channels
of particular interest.

\FloatBarrier
\begin{figure}[ht]
\captionsetup{justification=justified,
              singlelinecheck=false}
\includegraphics[width=\linewidth]
    {figures/solvency_stablecoins.png}
\caption{Diagnostic plots for the selected model --- solvency risk
$b_{t,j}$ (DAI and USDC, $s = 100\%$, 11 variables,
Train $R^2 = 0.958$, Test $R^2 = 0.929$).\\
\scriptsize{Note: 
\textit{Upper-left}: actual against fitted values on the training
set (blue).
\textit{Upper-right}: actual against fitted values on the
held-out test set (coral).
\textit{Lower-left}: HAC-estimated coefficients with
95\,\% confidence intervals, ordered by magnitude.
\textit{Lower-right}: residuals over time for the training period (blue) and the test period (coral), separated by the dashed
vertical line at the train/test boundary (March 2023).
}}
\label{fig:solvency_stablecoins_size100}
\end{figure}

The most prominent channel operates through \emph{own-pool financial characteristics}. Collateral utilization ($+$) and the collateral factor ($+$) carry the largest own-pool
coefficients, indicating that pools accepting more collateral at higher loan-to-value ratios generate more bad debt when the stablecoin de-pegs. The self-borrowing rate also enters positively, consistent with internally recycled positions
being more exposed to a sudden loss of parity.

The second channel operates through \emph{collateral network features}. The value-weighted out-weight HHI of collateral providers enters significantly, indicating that more concentrated collateral structures amplify bad debt under a de-pegging event, mirroring the finding from the crypto solvency analysis.

The third and distinctive channel, absent in the crypto models, operates through \emph{borrower characteristics}. Three borrower variables are retained by LASSO: the
value-weighted in-weight HHI of borrowers, their value-weighted out-degree, and their value-weighted collateral utilization rate. \latestrev{These selections indicate that borrower characteristics contain predictive information alongside collateral-provider characteristics. They do not separately identify how borrowers adjust their positions following a price decline.}

% ────────────────────────────────────────────────────────────────
\subsection{Incremental Value of Liability-Side Information}
\label{sec:results_incremental}

A key question motivating our empirical design is whether 
reconstructing the liability network from on-chain data adds 
predictive content beyond what can be inferred from the asset 
side of the pools' balance sheets as already published by 
Compound itself. To answer this directly, we estimate a sequence 
of nested models defined in \autoref{tab:incremental}, expanding 
the candidate feature set progressively from a single asset-side 
predictor to the full liability-side specification, \rev{ while \autoref{fig:incremental} displays the corresponding $R^2$ 
profiles. 

The test $R^2$ allows one to compare predictive performance
across models since it summarizes average accuracy over the whole
test panel. However, average accuracy is not the metric that
matters most for risk monitoring: a supervisor is less interested
in predicting the level of stress on every pool-day than in
identifying the few that will be severely stressed. We therefore
report, alongside test $R^2$, the top-decile precision of each
specification: we rank test pool-days by predicted risk, flag the
worst decile, and measure the share of flagged pool-days that are
genuinely in the worst decile of \latestrev{simulated} risk.\footnote{The test
panel contains $N_{\text{test}} = 401$ pool-days, so the flagged
decile comprises the $40$ pool-days with the highest predicted
value. Because the crypto sample contains only two pools per date,
pool-days are ranked across the test panel rather than within
date.} A ranking that carries no information attains $0.10$ in
expectation, and a perfect ranking attains $1$.
\ref{app:decision_metrics} reports three further metrics that bear
directly on risk monitoring. They yield the same
qualitative conclusion, with specifications that incorporate
topological information dominating the asset-side baseline on
every dimension.}


\FloatBarrier 
\begin{table}[ht]
\centering
\begin{threeparttable}
\caption{Incremental predictive value of liability-side
         information. }
\label{tab:incremental}
\footnotesize
\begin{tabular}{lrrrrrr}
\toprule
& \multicolumn{3}{c}{Liquidity risk ($\ell_{t,j}$)}
& \multicolumn{3}{c}{Solvency risk (\rev{$b_{t,j}$})} \\
\cmidrule(lr){2-4}\cmidrule(lr){5-7}
Model
    & Test $R^2$
    & \rev{Top-10\%}
    & Vars
    & Test $R^2$
    & \rev{Top-10\%}
    & Vars \\
\midrule
M1: Asset side
    & $0.236$ & \rev{$0.025$} & 1
    & $0.393$\rev{\textsuperscript{?}} & \rev{$0.050$} & 3\rev{\textsuperscript{?}} \\
M2: $+$ Pool features
    & $0.523$ & \rev{$0.400$} & 11
    & $0.731$ & \rev{$0.625$} & 8  \\
M3: $+$ Collateral network
    & $0.869$ & \rev{$0.600$} & 10
    & $0.749$ & \rev{$0.650$} & 11 \\
\textbf{M4: $+$ User features (full model)}
    & $\mathbf{0.868}$ & \rev{$\mathbf{0.600}$} & $\mathbf{14}$
    & $\mathbf{0.834}$ & \rev{$\mathbf{0.750}$} & $\mathbf{10}$ \\
\bottomrule
\end{tabular}
\begin{tablenotes}[flushleft]
\setlength{\itemsep}{0pt}
\setlength{\parsep}{0pt}
{\scriptsize
\item \textit{Note}: M1: collateral factor and utilization
ratio. M2: $+$ pool financial characteristics, centrality,
and cross-pool linkages. M3: $+$ exposure-weighted
connectivity of collateral providers. M4: $+$ financial
condition of collateral providers (health ratio, collateral
utilization, self-borrowing rate, deposit share).
\rev{Top-10\%: top-decile precision, defined in the text.}
LASSO
penalty selected by expanding-window cross-validation
($K = 5$ folds); OLS refit on selected variables with
Newey--West HAC errors ($L = 4$ lags); trimming retains
predictors significant at $\alpha = 0.001$. Predictors
standardized using training-set moments only.
\rev{Equality of predictive accuracy between successive rows is
assessed with the one-sided \citet{clark2007approximately}
statistic, computed on the nested specifications estimated without
variable selection so that the smaller model is exactly nested in
the larger one as the test requires; the Test $R^2$ and Top-10\%
columns report the selected and trimmed specifications used
throughout the paper. The direct comparison of M1 against M4
rejects sufficiency of the asset-side information set at
$p < 0.001$ for both outcomes.}
}
\end{tablenotes}
\end{threeparttable}
\end{table}

\FloatBarrier 
\begin{figure}[ht]
\captionsetup{justification=justified, singlelinecheck=false}
\begin{subfigure}[b]{0.49\textwidth}
    \centering
    \includegraphics[width=\textwidth]{figures/incremental_liquidity.png}
\end{subfigure}
\hfill
\begin{subfigure}[b]{0.49\textwidth}
    \centering
    \includegraphics[width=\textwidth]{figures/solvency_risk_increment.png}
\end{subfigure}
\caption{Incremental predictive value of liability-side
information (ETH and WBTC, $s = 50\%$).\\
\scriptsize{Note:
M1 includes the collateral factor and utilization ratio as
asset-side predictors. M2 adds pool financial
characteristics, centrality measures, and cross-pool
linkages. M3 adds exposure-weighted connectivity measures
of collateral providers, constructed from the bipartite
liability network. M4 adds the financial condition of
collateral providers (health ratio, collateral utilization,
self-borrowing rate, deposit share), completing the full
model. Solid circles show test $R^2$; dashed squares show
train $R^2$.
\textit{Left panel}: liquidity risk $\ell_{t,j}$.
\textit{Right panel}: solvency risk $b_{t,j}$.}}
\label{fig:incremental}
\end{figure}
\FloatBarrier 

\boldpara{Liquidity risk.}
The asset-side baseline (M1), which includes only the collateral 
factor and utilization ratio, explains $23.6\%$ of out-of-sample variation in liquidity 
risk ($\ell_{t,j}$). Adding the full set of pool features --- 
financial characteristics, centrality measures, and cross-pool 
linkages --- at M2 raises test $R^2$ to $0.523$ ($+28.7$ 
percentage points). The dominant gain in the sequence comes at 
M3, where adding collateral network features raises test $R^2$ 
to $0.869$ ($+34.6$ percentage points). Adding the financial 
condition of depositors at M4 yields no further improvement in 
out-of-sample fit ($R^2 = 0.868$). \rev{A Clark--West test on the 
exactly nested specifications does not reject equal predictive 
accuracy between M3 and M4 \latestrev{at the 5\% level} (one-sided 
$p = 0.071$), whereas each earlier step rejects at $p < 0.001$.} 
\latestrev{In the selected models, depositor financial characteristics add little to held-out fit once connectivity is included.} This pattern confirms that, for reserve 
depletion, \emph{connectivity is the primary channel}: knowing 
how broadly depositors are linked across the protocol is 
sufficient to predict liquidation-driven liquidity losses, and 
knowing how leveraged they are adds little beyond that.

\medskip

\boldpara{Solvency risk.}
The pattern is less stark for solvency risk ($b_{t,j}$). The
asset-side baseline (M1) explains $39.3\%$ of out-of-sample
variation. Adding pool features at M2 raises test $R^2$ to
$0.731$ ($+33.8$ percentage points) --- a larger gain than in
the liquidity case, reflecting that pool-level network position
(particularly PageRank centrality) is an important predictor of
bad debt. Adding the collateral network at M3 delivers only a
small gain to $0.749$ ($+1.8$ percentage points), \rev{which the
same test does not distinguish from zero (one-sided $p = 0.96$).}
Adding the financial condition of collateral providers at M4
yields a final improvement to $0.834$ ($+8.5$ percentage points)
\rev{and rejects at $p < 0.001$: for bad debt it is therefore the
balance-sheet condition of collateral providers, rather than their
connectivity, that carries the incremental information}.
The financial condition of collateral providers thus contributes
more at the final step than their network connectivity, the
reverse of the pattern for liquidity risk. This reflects the
different transmission mechanism: bad debt arises when collateral
\emph{value} fails to cover outstanding debt, a condition that
depends jointly on the network position of depositors --- which
determines the scale of liquidation cascades --- and on their
financial condition, which determines whether individual positions
generate unrecoverable shortfalls.

\medskip
\boldpara{The asset-side gap.}
Taken together, these results establish that the information
published directly by Compound, namely the collateral factor
and utilization ratio that constitute our asset-side baseline,
provides only a partial basis for risk monitoring. A surveillance
system relying solely on these observable pool-level indicators
would explain less than a quarter of the variation in liquidity
risk ($R^2 = 0.236$) and less than $40\%$ of the variation in
solvency risk ($R^2 = 0.393$). \rev{The ranking such a system
induces is no better than chance: of the pool-days it flags as
most at risk, $2.5\%$ are genuinely in the worst decile of
\latestrev{simulated} reserve depletion, against the $10\%$ a blind ranking
would attain, and $5.0\%$ for bad debt against the same benchmark.
By contrast, the full model reaches $60\%$ and $75\%$ respectively.} Even
augmenting the asset side
with the full set of pool financial characteristics, centrality
measures, and cross-pool linkages (M2) leaves substantial
predictive gaps of $34.5$ percentage points for liquidity risk
and $10.3$ percentage points for solvency risk relative to the
full model. These findings underscore that topology-aware monitoring is not a mere refinement of standard risk dashboards but a necessary input for assessing systemic vulnerability in decentralized lending protocols.

%Each row adds one information set to the candidate pool. LASSO selects variables within the cumulative candidate set; OLS is refit on selected variables with Newey--West HAC standard errors ($L = 4$ lags); a final trimming step  retains only predictors significant at  $\alpha = 0.001$. $\Delta$Test $R^2$ is the gain   relative to the previous row. ETH and WBTC, $s = 0.50$.


\section{\rev{External Validity}}
\label{sec:external_validity}

\rev{One may wonder whether our findings carry over to more recent
lending protocols, since Compound V2 is no longer the largest platform and its
activity declined markedly after the 2021--2022 peak. However, this decline does not make the insights drawn
from our study obsolete, because Compound V2 provided the template for
the protocols that succeeded it. \latestrev{Collateralized borrowing and rule-based liquidation remain important features of later protocols, but the treatment of collateral, shared liquidity, and losses differs.} \autoref{tab:protocol_designs} sets out the main differences
between Compound V2 and its successors, \latestrev{with a fuller, sourced comparison in \ref{app:later_protocols}}.}

\FloatBarrier
\begin{table}[ht]
\centering
\begin{threeparttable}
\caption{\rev{Design of the main pooled lending protocols.}}
\label{tab:protocol_designs}
\footnotesize
\begin{tabularx}{\textwidth}{l >{\raggedright\arraybackslash}X >{\raggedright\arraybackslash}X >{\raggedright\arraybackslash}X}
\toprule
\rev{Protocol} & \rev{Market structure} & \rev{Liquidation rule}
& \rev{\latestrev{Collateral and liquidity links}} \\
\midrule
\rev{Compound V2} &
\rev{\latestrev{Shared pools; eligible collateral backs permitted borrowing.}} &
\rev{\latestrev{Close factor: up to $50\%$ of each borrowed asset per transaction.}} &
\rev{\latestrev{User accounts link collateral and borrowing pools.}} \\[4pt]
\rev{Compound III} &
\rev{\latestrev{One borrowable base asset per deployment; other collateral is not lent out.}} &
\rev{Whole-account absorption; collateral resold at a discount.} &
\rev{\latestrev{Collateral backs base-asset debt; no V2-style collateral lending pool.}} \\[4pt]
\rev{Aave V3} &
\rev{\latestrev{Shared reserves with version-specific collateral and borrowing restrictions.}} &
\rev{\latestrev{Close-factor limits vary with position health, size, and deployed version.}} &
\rev{\latestrev{Permitted collateral--debt links within a market; restrictions vary by version.}}
\\[4pt]
\rev{Aave V4} &
\rev{Shared liquidity hub with market-specific spokes.} &
\rev{Liquidation to a target health factor, with a bonus scaling in
the shortfall.} &
\rev{\latestrev{Collateral accounts are spoke-specific; hubs supply shared liquidity.}} \\[4pt]
\rev{\latestrev{Morpho Blue}} &
\rev{Isolated markets, one collateral asset and one loan asset each.} &
\rev{Fixed incentive set by the market's liquidation loan-to-value.} &
\rev{\latestrev{Market-level positions; vaults may hold claims across markets.}} \\
\bottomrule
\end{tabularx}
\begin{tablenotes}[flushleft]
\setlength{\itemsep}{0pt}
\setlength{\parsep}{0pt}
{\scriptsize
\item \rev{\textit{Note}: Compound V2 is the protocol studied in
this paper. \latestrev{Collateral links, shared funding, and shared loss exposure are distinct. The descriptions of later protocols are based on the primary sources and version qualifications in \ref{app:later_protocols}.}}
}
\end{tablenotes}
\end{threeparttable}
\end{table}
\FloatBarrier

\rev{\autoref{tab:protocol_designs} shows that innovation has
proceeded along two dimensions. On the one hand, protocols have
refined the liquidation mechanism, \latestrev{in some cases} moving away from the discrete close
factor of Compound V2 towards rules that calibrate the amount seized
to the size of the shortfall. On the other, they have sought to
mitigate the channels through which contagion spreads, by segregating
collateral, capping exposures, or isolating markets altogether. We
now consider each dimension in turn.}

\medskip

\label{sec:liquidation_rule}
\rev{\boldpara{The liquidation rule.}
Our propagation algorithm liquidates the amount that returns a
borrower's health factor to one (Assumption \ref{hyp:H1}). \latestrev{Aave V4 also uses a target health factor, but its configured target, bonus schedule, and treatment of small residual positions differ from our assumptions (\ref{app:later_protocols_aave4}). Compound V2 instead caps repayment of an individual borrowed asset in each liquidation transaction. Simulating health-restoring and close-factor liquidation on the same Compound V2 positions therefore assesses both our approximation and the mechanical effect of changing this rule; it does not estimate the overall effect of Aave V4's implementation.}}


\rev{Panel A of \autoref{tab:liquidation_rule} reports mean reserve
depletion across crypto pool-days under each rule. The close factor
destroys more reserves at every shock size below $70\%$, and the
excess declines with the size of the shock: $59.3\%$ at a $20\%$ shock, $26.6\%$ at
$30\%$, $15.3\%$ at $40\%$, $7.5\%$ at our $50\%$ benchmark, and
$3.9\%$ at $60\%$. This pattern follows from the structure of the
rule. \latestrev{Repayment up to the cap can exceed the amount needed to restore health}, which is the case for
positions that have only just crossed the liquidation threshold. Once
a shock is severe enough that positions are deeply undercollateralized,
the health-restoring amount approaches the full collateral balance and
the two rules converge. Equivalently, the excess can be used to gauge
the error our algorithm makes in representing Compound V2: \latestrev{the close-factor outcome is $7.5\%$ higher at the benchmark shock, with a gap below $4\%$ for larger shocks}. Furthermore, at the benchmark shock,
the close factor leaves the ranking of
pools essentially unchanged, with a rank correlation of $0.98$ between
the two outcome series. \latestrev{At the benchmark shock, the close-factor rule therefore raises mean reserve depletion modestly while largely preserving the ranking of pool-days.}}

\rev{Panel B shows that predictive accuracy does not depend on the
liquidation rule. The full model attains a test $R^2$ of $0.882$ for
liquidity risk when outcomes are generated under the close factor,
against $0.868$ under health-restoring liquidation, and $0.849$ for
solvency risk under either rule. Bad debt is identical under the two rules \latestrev{with the other simulation assumptions held fixed}: bad debt arises when collateral value fails to cover debt
after the liquidation bonus, a condition that does not depend on the
order or the size of individual seizures.}

\rev{\latestrev{To summarize, health-restoring liquidation reduces the amount seized from positions that have only just crossed the liquidation threshold. The difference is smaller for the severe shocks considered here, while the ranking of pool-days and the predictive results are similar. These comparisons quantify the mechanical consequences of the rule on fixed positions, rather than borrowers' or liquidators' responses to adopting it.}}

\FloatBarrier
\begin{table}[ht]
\centering
\begin{threeparttable}
\caption{\rev{Health-restoring versus close-factor liquidation.}}
\label{tab:liquidation_rule}
\footnotesize

\begin{tabular*}{0.85\textwidth}{l@{\extracolsep{\fill}}rrr}
\multicolumn{4}{c}{\rev{\textit{Panel A. Mean reserve depletion, crypto
pools, percent of pool assets.}}} \\[1pt]
\toprule
\rev{Shock $s$} & \rev{Health-restoring} & \rev{Close factor}
& \rev{Excess} \\
\midrule
\rev{$0.20$} & \rev{$0.84$}  & \rev{$1.33$}  & \rev{$+59.3\%$} \\
\rev{$0.30$} & \rev{$3.68$}  & \rev{$4.66$}  & \rev{$+26.6\%$} \\
\rev{$0.40$} & \rev{$7.23$}  & \rev{$8.33$}  & \rev{$+15.3\%$} \\
\rev{$\mathbf{0.50}$} & \rev{$\mathbf{11.17}$} & \rev{$\mathbf{12.01}$}
& \rev{$\mathbf{+7.5\%}$} \\
\rev{$0.60$} & \rev{$13.87$} & \rev{$14.41$} & \rev{$+3.9\%$} \\
\rev{$0.70$} & \rev{$14.31$} & \rev{$14.28$} & \rev{$-0.2\%$} \\
\bottomrule
\end{tabular*}

\vspace{4pt}

\begin{tabular*}{0.85\textwidth}{l@{\extracolsep{\fill}}lrr}
\multicolumn{4}{c}{\rev{\textit{Panel B. Out-of-sample predictive
performance.}}} \\[1pt]
\toprule
& & \multicolumn{2}{c}{\rev{Test $R^2$}} \\
\cmidrule(lr){3-4}
\rev{Outcome} & \rev{Model} & \rev{Health-restoring}
& \rev{Close factor} \\
\midrule
\rev{Liquidity} & \rev{M1: Asset side}
    & \rev{$0.236$} & \rev{$0.191$} \\
\rev{Liquidity} & \rev{M4: Full model, frozen}
    & \rev{$0.868$} & \rev{$0.882$} \\
\rev{Solvency}  & \rev{M1: Asset side}
    & \rev{$0.397$\textsuperscript{?}} & \rev{$0.397$\textsuperscript{?}} \\
\rev{Solvency}  & \rev{M4: Full model, frozen $+$ FE}
    & \rev{$0.849$} & \rev{$0.849$} \\
\bottomrule
\end{tabular*}

\begin{tablenotes}[flushleft]
\setlength{\itemsep}{0pt}
\setlength{\parsep}{0pt}
{\scriptsize
\item \rev{\textit{Note}: Both rules are simulated on the same daily
network snapshots. Panel A omits $s = 0.10$, where both levels are
near zero, and $s \geq 0.80$, where both rules liquidate positions in
full. Panel B: M1 and M4 are the specifications of
\autoref{tab:incremental}; the M4 rows are held fixed at the variables
selected under the health-restoring rule and refitted on each outcome, so the two columns differ only in the
dependent variable. The solvency rows include the WBTC fixed effect
and are reported before post-HAC re-trimming, which gives the $0.834$
of \autoref{tab:incremental}.}
}
\end{tablenotes}
\end{threeparttable}
\end{table}
\FloatBarrier

\medskip

\label{sec:market_structure}
\rev{\boldpara{Market structure.}
The other main stream of updates has restricted the links
that users may create across pools, with the explicit aim of
containing contagion. The
designs recorded in \autoref{tab:protocol_designs} follow from this
concern. Compound III confines borrowing to a single base asset and
segregates collateral, so that a user cannot connect two borrowing
pools within a deployment. Aave V3 retains shared pools but \latestrev{limits collateral and borrowing combinations through version-specific controls, including the original isolation mode and later eMode configurations}. Morpho \latestrev{Blue uses separate lending markets}, pairing one
collateral asset with one loan asset in each market. \latestrev{Vaults can hold claims across several such markets without automatically transferring one market's bad debt to another; \ref{app:later_protocols} explains this distinction.}}

\rev{We cannot devise a counterfactual to evaluate these changes as we did 
for the liquidation
rule. Whereas the close factor changes only how a given network is
unwound, restrictions on which pools a user may connect would affect
the formation of the network itself. Borrowers and lenders would redistribute
their positions across markets, and the resulting topology would
differ fundamentally from the one we observe. A counterfactual holding
positions fixed would ignore these endogenous responses, yielding both
unfeasible network configurations under the new constraints and
unreliable estimates.}

\rev{Our results nonetheless give empirical grounding to the concern
these new designs address. According to our findings, the exposure-weighted connectivity of
collateral providers is the single largest source of predictive
content for reserve depletion (see \autoref{tab:incremental}). Restricting
cross-pool links is therefore aimed at a channel that we show to be
quantitatively important, rather than at a hypothetical one. A natural
direction for further research would be to apply our methodology to these
new networks, and to measure how far the
restrictions have in fact contained systemic risk.}


\section{Conclusion}
\label{sec:conclusion}

This paper provides a high-frequency, network-based empirical analysis of systemic risk in decentralized lending protocols using a full reconstruction of the Compound protocol's liability network from 2020 to 2024. After computing several topological and financial measures, we quantify the effects of interdependencies and find that \rev{the indicators the protocol publishes capture only a fraction of the stress it experiences, whereas measures built from the liability network predict most of it}. These findings suggest that systemic risk in DeFi is not a function of size or collateral coverage alone, but of how users and assets are embedded in the network of obligations.

Our findings speak to a broader opportunity created by DeFi's architecture. Unlike traditional financial systems, where network reconstruction depends on confidential exposures and reporting lags, DeFi protocols record every transaction on a public blockchain in real time, allowing balance sheets and liability links to be mapped continuously. Furthermore, smart contracts encode the system's responses to stress through liquidation rules and interest rate adjustments, making it possible to simulate the propagation of shocks with a precision that is unattainable in traditional finance. Together, these features transform the feasibility of systemic risk assessment: rather than relying on periodic snapshots of aggregated balance sheet indicators, regulators and protocol governance bodies could monitor the full network of obligations and simulate system responses in real time. Our methodology is designed to exploit this opportunity, with results suggesting that doing so captures dimensions of systemic vulnerability that aggregate indicators cannot reveal. As policymakers and central banks consider how new infrastructures---such as CBDCs, Regulated Liability Networks, and tokenized capital markets---might interact with DeFi ecosystems, our findings indicate that understanding network externalities will be critical to evaluating financial resilience.

Several directions remain open for future research. Our simulations could be enriched by incorporating more complex price dynamics, including gradual price declines and intra-day volatility, as well as user-level behavioral responses such as collateral refilling and endogenous deleveraging. Extending the analysis to block-level data would allow the model to be calibrated against intra-day liquidation dynamics, opening the door to very short-term predictions relevant to traders and arbitrageurs rather than the regulatory and governance audience targeted here. Finally, our framework focuses on liquidation-driven propagation but does not model endogenous liquidity runs, in which depositors withdraw en masse, pushing utilization ratios toward full depletion and trapping lenders. The recent liquidity crunch in Aave triggered by the KelpDAO exploit of April 2026 shows that such risk is a practical concern.\footnote{A bridge vulnerability in KelpDAO's infrastructure enabled the minting of unbacked rsETH tokens, which were used as collateral on Aave to borrow ETH. Once the hack became public, it triggered cascading withdrawals across Aave's ETH, USDT and USDC pools, driving utilization to 100\% and generating substantial bad debt, with contagion spreading from directly exposed ETH lenders to stablecoin lenders \citep{bristow2026}. See also \href{https://governance.aave.com/t/rseth-incident-report-april-20-2026/24580}{Aave's incident report} (https://governance.aave.com/t/rseth-incident-report-april-20-2026/24580) for further details.} Extending our balance sheet reconstruction and propagation framework to incorporate depositor run dynamics would therefore be a natural and policy-relevant next step.


\bibliography{refs}

\newpage

\appendix

\section{Aggregation Procedure}
\label{app:Aggregation}

We explain in this appendix how to reconstruct the pool's balance sheet after liquidations. We denote by \(\mathbf{C}_j\) and \(\mathbf{D}_j\) the amounts borrowed and deposited in pool j, so that:
\[
\mathbf{C}_j = \sum_{i=1}^M L_{ij}, \quad \mathbf{D}_j = \sum_{i=1}^M L_{ji},
\]
where $M = |\mathcal{U}|$ is the number of users. We also use \(\mathbf{T}_j\) to denote the available tokens in pool \(j\) (i.e., deposited but not borrowed), and \(\mathbf{E}_j\) the equity of pool \(j\). Considering that pools have no bad debt initially, the balance sheet of pool \(j\) is given by
\begin{align*}
\begin{cases}
    \text{Assets}_j \quad&=\quad \mathbf{T}_j + \mathbf{C}_j, \\
    \text{Liabilities}_j\quad &=\quad \mathbf{D}_j + \mathbf{E}_j.
\end{cases}
\end{align*}

We now characterize the impact of a price shock. We denote by \(\mathcal{L}\) and \(\mathcal{L}'\) the subsets of liquidated users without and with bad debt after the shock, respectively. As explained in \autoref{sec:Propagation_Algo}, the liquidated amounts $\Delta_i$ for each user are given by
\begin{equation*}
    \Delta_i =
    \begin{cases}
        \dfrac{D_i - CL_i}{D_i \cdot \left[1 - (1+LB)\cdot \dfrac{CL_i}{C_i}\right]}, & \text{when } i \in \mathcal{L},\\[6pt]
        \dfrac{C_i}{(1+LB)\cdot D_i}, & \text{when } i \in \mathcal{L}',
    \end{cases}
\end{equation*}
where \(C_i\), \(D_i\), and $CL_i$ denote the collateral, debt, and pledgeable collateral of user \(i\), as defined in \eqref{eq:CiDi} and \eqref{eq:CLi} respectively. Aggregating their values across all users yields the following updating rules for the components of the pools' balance sheets:
\begin{align*}
    \mathbf{T}_j' &= \mathbf{T}_j + \sum_{i\in \mathcal{L} \cup \mathcal{L}'} \Delta_i \cdot L_{ij} - (1 + LB)\cdot\sum_{i\in \mathcal{L} \cup \mathcal{L}'} \Delta_i\cdot \frac{D_i}{C_i}\cdot L_{ji},\\
    \mathbf{C}_j' &= \mathbf{C}_j - \sum_{i\in \mathcal{L}'}L_{ij} - \sum_{i\in \mathcal{L}} \Delta_i \cdot L_{ij},\\
    \mathbf{D}_j' &= \mathbf{D}_j - (1 + LB)\cdot \sum_{i\in \mathcal{L}\cup\mathcal{L}'} \Delta_i \cdot \frac{D_i}{C_i}\cdot L_{ji},\\
    \mathbf{B}_j' &= \sum_{i\in \mathcal{L}'}  (1 - \Delta_i)\cdot L_{ij},
\end{align*}
where primes denote post-shock values and $\mathbf{B}_j$ is the bad debt in pool $j$. By definition, we have \(\mathbf{E}_j' = \mathbf{E}_j - \mathbf{B}_j'\). Using \(k\) to index the shocked pool, the total bad debt generated by the price shock to \(k\) is given by
\begin{equation}
\label{eq:Def_BD}
    BD_k = \sum_{j=1}^p \mathbf{B}'_j,
\end{equation}
where $p=|\mathcal{P}|$ is the number of pools. When liquidation volumes are large enough to deplete the reserves of a pool, the equations above must be adjusted to account for the inability of liquidators to redeem their \texttt{cTokens}. Since unredeemed \texttt{cTokens} remain in the pools, they are added to the post-liquidation deposits:
\begin{align*}
    \mathbf{T}_j' &= \max\left\{0,\mathbf{T}_j + \sum_{i\in \mathcal{L} \cup \mathcal{L}'} \Delta_i \cdot L_{ij} - (1 + LB)\cdot\sum_{i\in \mathcal{L} \cup \mathcal{L}'} \Delta_i\cdot \frac{D_i}{C_i}\cdot L_{ji}\right\},\\
    \mathbf{D}_j' &= \mathbf{D}_j - (1 + LB)\cdot \sum_{i\in \mathcal{L}\cup\mathcal{L}'} \Delta_i \cdot \frac{D_i}{C_i}\cdot L_{ji}\\
    &+\max\left\{0,(1 + LB)\cdot\sum_{i\in \mathcal{L} \cup \mathcal{L}'} \Delta_i\cdot \frac{D_i}{C_i}\cdot L_{ji}-\mathbf{T}_j - \sum_{i\in \mathcal{L} \cup \mathcal{L}'} \Delta_i \cdot L_{ij} \right\}.\\
\end{align*} 

\section{Feature Extraction}\label{app:feat_des_app}

The bipartite network represents liability relationships between user and lending pool nodes, denoted as $L_{ij}$, where node $i$ has a liability to node $j$. Based on the daily asset prices, the edge weights correspond to the liability amounts converted to U.S. dollars.

\textbf{Node features.} We extract node-level features as summarized in \autoref{tab:node_features}. We further compute the share of each lending pool's equity, external assets, and size to assess the relative scale of each lending pool. These shares are calculated by normalizing each variable by the total across all pool nodes, i.e., $share_j = variable_j / \sum_k variable_k$, where $j$ indexes a specific lending pool and $k$ ranges over all pools.

\FloatBarrier 
\begin{table}[ht]
\caption{Node features for users and lending pools}
\label{tab:node_features}
\resizebox{\textwidth}{!}{%
\begin{tabular}{lll}
\toprule
\textbf{Feature} & \textbf{User Node} & \textbf{Lending Pool Node} \\
\midrule
In-degree         & Number of deposited assets       & Number of borrowing users        \\
In-weight         & Total deposited amount           & Total borrowed amount            \\
In-weight HHI     & Herfindahl index of deposited assets & Herfindahl index of borrowing users \\
\rowcolor[HTML]{F3F3F3}
Out-degree        & Number of borrowed assets        & Number of depositing users       \\
\rowcolor[HTML]{F3F3F3}
Out-weight        & Total borrowed amount            & Total deposited amount           \\
\rowcolor[HTML]{F3F3F3}
Out-weight HHI    & Herfindahl index of borrowed assets & Herfindahl index of depositing users \\
In/Out Ratio      & Deposit-to-borrow ratio          & Borrow-to-deposit ratio          \\
Health Ratio      & Health ratio of users            & N/A                              \\
External assets   & 0 (assuming no external assets)  & Available liquidity (cash) + reserves \\
\rowcolor[HTML]{F3F3F3}
Equity            & \multicolumn{2}{c}{External assets + In-weight -- Out-weight}         \\
\rowcolor[HTML]{F3F3F3}
Size              & \multicolumn{2}{c}{In-weight + Out-weight}                           \\
\rowcolor[HTML]{F3F3F3}
\% Self-borrowing    & \multicolumn{2}{c}{Percentage of self-borrowed amount} \\
\bottomrule
\end{tabular}}
\begin{minipage}{\textwidth}
\footnotesize \textit{Note:} The Herfindahl-Hirschman Index (HHI) is calculated as $HHI = \sum_i s_i^2$, where $s_i$ is the share of value associated with asset $i$ (for users: deposits or borrows; for pools: users depositing or borrowing). A higher HHI indicates greater concentration.
\end{minipage}
\end{table}

In addition to these attributes, we compute several centrality measures for each node to characterize its structural role in the network:
\begin{itemize}
    \item \emph{Closeness Centrality}, estimated using approximate algorithms provided by the NetworKit library~\citep{angriman2023algorithms}, due to computational limitations of exact methods;
    \item \emph{Hubs and Authorities}, computed via the HITS algorithm~\citep{kleinberg1999authoritative};
    \item \emph{PageRank}, which we use as a practical proxy for eigenvector centrality, given its faster convergence and scalability~\citep{page1999pagerank}.
\end{itemize}

\textbf{Edge features.} For each user–pool edge, we compute the share of deposits and borrows to quantify the user’s relative exposure to the pool. These features are defined as:
\begin{equation}
\text{shareDeposit}_{u,i} = \frac{\text{deposit}_{u,i}}{\sum_{i} \text{deposit}_{u,i}}, \qquad
\text{shareBorrow}_{u,i} = \frac{\text{borrow}_{u,i}}{\sum_{i} \text{borrow}_{u,i}},
\end{equation}
\noindent where $\text{deposit}_{u,i}$ and $\text{borrow}_{u,i}$ denote the amounts deposited to or borrowed from pool $i$ by user $u$. These ratios capture the proportion of a user's activity attributed to a specific pool.


\textbf{Lenders and borrowers features:} We compute the weighted average of node attributes for lenders (nodes with incoming edges from the pool) and borrowers (nodes with outgoing edges to the pool), using the out-weight and in-weight, respectively, as weighting factors. This allows us to capture the average characteristics of users interacting with each pool, adjusted by the magnitude of their deposited or borrowed amounts.

\autoref{tab:regression_features} lists the features included in the regression models. 
The features include topological metrics and financial attributes for the shocked pool, as well as user-aggregated statistics for lenders and borrowers. 


\FloatBarrier 
\begin{table}[ht]
\centering
\begin{threeparttable}
\caption{Predictor variables used in
         the regression models for
         liquidity risk ($\ell_{t,j}$)
         and solvency risk ($b_{t,j}$).
         Variables are organized by
         predictor group. Collateral
         provider and borrower statistics
         are value-weighted averages
         computed at the pool level.}
\label{tab:regression_features}
\footnotesize
\begin{tabular}{lccccc}
\toprule
& \multicolumn{2}{c}{Pool level}
& \multicolumn{2}{c}{Collateral (wavg.)}
& \multicolumn{1}{c}{Borrowers (wavg.)} \\
\cmidrule(lr){2-3}
\cmidrule(lr){4-5}
\cmidrule(lr){6-6}
Variable
    & Own-pool
    & Cross-pool
    & Network
    & Financial
    & Network \& Financial \\
\midrule
\multicolumn{6}{l}{%
    \textit{Topological features}} \\
\quad In-degree
    & \checkmark & & \checkmark & & \checkmark \\
\quad Out-degree
    & \checkmark & & \checkmark & & \checkmark \\
\quad In-weight HHI
    & \checkmark & & \checkmark & & \checkmark \\
\quad Out-weight HHI
    & \checkmark & & \checkmark & & \checkmark \\
\quad In-flow weight
    & & & \checkmark & & \checkmark \\
\quad In-flow weight HHI
    & & & \checkmark & & \checkmark \\
\quad Out-flow weight HHI
    & & & \checkmark & & \checkmark \\
\quad PageRank
    & \checkmark & & & & \\
\quad Hub score
    & \checkmark & & & & \\
\quad In-degree dispersion
    & & \checkmark & & & \\
\quad In-degree (wavg.)
    & & \checkmark & & & \\
\quad Out-degree (wavg.)
    & & \checkmark & & & \\
\quad Deposit concentration (HHI)
    & & \checkmark & & & \\
\quad Borrowing concentration (HHI)
    & & \checkmark & & & \\
\quad Self-borrowing dispersion
    & & \checkmark & & & \\
\quad Self-borrowing share (wavg.)
    & & \checkmark & & & \\
\addlinespace
\multicolumn{6}{l}{%
    \textit{Financial features}} \\
\quad Collateral utilization
    & \checkmark & & & \checkmark & \checkmark \\
\quad Utilization ratio
    & \checkmark & & & & \\
\quad Self-borrowing rate
    & \checkmark & & & \checkmark & \checkmark \\
\quad Collateral factor
    & \checkmark & & & & \\
\quad Health ratio
    & & & & \checkmark & \checkmark \\
\quad Deposit share
    & & & & \checkmark & \checkmark \\
\bottomrule
\end{tabular}
\begin{tablenotes}[flushleft]
\setlength{\itemsep}{0pt}
\setlength{\parsep}{0pt}
\scriptsize
\item \textit{Notes}: HHI: Herfindahl--Hirschman index.
Wavg.: value-weighted average over collateral providers
or borrowers at the pool level. Cross-pool variables are
protocol-level and identical across pools on any given
date. ``Self-borrowing rate'' refers to distinct
variables at the pool, collateral provider, and borrower
levels. Borrower features are included in the stablecoin
specification and in the robustness checks; they are excluded from
the baseline crypto specification.
\end{tablenotes}
\end{threeparttable}
\end{table}

%\newpage
%%%%%%%%%%%%%%%%%%%%%%%%%%%%%%%%%%%%%%%%%%%%%%%%%%%%%%%%%%%%%%%%%%
%%%%% New tables
\FloatBarrier
\section{Descriptive Statistics}
\FloatBarrier 
\begin{table}[H]
\centering
\begin{threeparttable}
\caption{Descriptive statistics (Panel A) ---
         Pool characteristics and collateral provider features,
         1 January 2021 -- 15 June 2024}
\label{tab:desc_stats_a}
\footnotesize
\begin{tabular}{lrrrrrrrr}
\toprule
& \multicolumn{2}{c}{ETH} & \multicolumn{2}{c}{WBTC}
& \multicolumn{2}{c}{DAI} & \multicolumn{2}{c}{USDC} \\
\cmidrule(lr){2-3}\cmidrule(lr){4-5}
\cmidrule(lr){6-7}\cmidrule(lr){8-9}
Variable & Mean & Std & Mean & Std & Mean & Std & Mean & Std \\
\midrule
\multicolumn{9}{l}{\textit{Outcome variables}} \\
Liquidity risk ($\ell$) & 0.1294 & 0.0413 & 0.0941 & 0.0449 & \multicolumn{2}{c}{---} & \multicolumn{2}{c}{---} \\
Bad debt share ($b$) & 0.0236 & 0.0151 & 0.0077 & 0.0062 & 0.0168 & 0.0192 & 0.1036 & 0.0696 \\
\midrule
\multicolumn{9}{l}{\textit{Own-pool characteristics}} \\
Collateral utilization & 0.3764 & 0.0393 & 0.3472 & 0.0441 & 0.3961 & 0.1841 & 0.3346 & 0.1131 \\
Self-borrowing rate & 0.0161 & 0.0155 & 0.0205 & 0.0286 & 0.3744 & 0.2061 & 0.2115 & 0.1778 \\
Utilization ratio & 0.0409 & 0.0222 & 0.0350 & 0.0334 & 0.5991 & 0.1523 & 0.6042 & 0.1610 \\
Collateral factor & 0.8043 & 0.0328 & 0.6869 & 0.0272 & 0.8073 & 0.0362 & 0.8196 & 0.0445 \\
In-degree & 1310.81 & 218.82 & 449.57 & 19.35 & 2786.53 & 136.81 & 2578.44 & 250.72 \\
In-weight HHI & 0.1988 & 0.0972 & 0.2577 & 0.1054 & 0.1782 & 0.0651 & 0.1034 & 0.0343 \\
Out-degree & 70102.56 & 4105.48 & 3310.23 & 218.33 & 18512.81 & 924.06 & 217109.73 & 998.96 \\
Out-weight HHI & 0.0508 & 0.0191 & 0.0703 & 0.0316 & 0.1290 & 0.0530 & 0.0682 & 0.0413 \\
\midrule
\multicolumn{9}{l}{\textit{Pool centrality}} \\
PageRank (scaled) & 0.4887 & 0.1033 & 0.1372 & 0.0132 & 0.9956 & 0.0168 & 0.9098 & 0.0577 \\
Hub score (scaled) & 0.4776 & 0.4219 & 0.3832 & 0.4244 & 0.3936 & 0.4561 & 0.2264 & 0.3530 \\
\midrule
\multicolumn{9}{l}{\textit{Cross-pool linkages}$^{a}$} \\
\cmidrule(lr){2-3}
\multicolumn{1}{l}{} & \multicolumn{2}{c}{All assets} & \multicolumn{6}{c}{} \\
\cmidrule(lr){2-3}
Deposit concentration (HHI) & 0.2188 & 0.0205 & \multicolumn{6}{c}{} \\
Borrowing concentration (HHI) & 0.3393 & 0.0296 & \multicolumn{6}{c}{} \\
In-degree dispersion & 901.00 & 79.34 & \multicolumn{6}{c}{} \\
Out-degree dispersion & 54425.86 & 4130.41 & \multicolumn{6}{c}{} \\
Self-borrowing dispersion & 0.2726 & 0.0836 & \multicolumn{6}{c}{} \\
Out-degree (wavg.) & 68577.29 & 14517.57 & \multicolumn{6}{c}{} \\
In-degree (wavg.) & 2309.51 & 273.74 & \multicolumn{6}{c}{} \\
Self-borrowing share (wavg.) & 0.3847 & 0.2249 & \multicolumn{6}{c}{} \\
\midrule
\multicolumn{9}{l}{\textit{Collateral network features}} \\
In-degree (wavg.) & 2.0017 & 0.4203 & 2.3651 & 0.3465 & 1.7121 & 0.5971 & 1.7050 & 0.4578 \\
In-flow weight (wavg.) & 160.95M & 217.55M & 122.86M & 128.69M & 256.37M & 276.89M & 193.90M & 247.64M \\
In-flow weight HHI (wavg.) & 0.8273 & 0.0391 & 0.7649 & 0.0538 & 0.9101 & 0.0964 & 0.9022 & 0.0685 \\
Out-degree (wavg.) & 1.7533 & 0.3142 & 1.9325 & 0.2936 & 1.4151 & 0.3986 & 1.4778 & 0.3929 \\
Out-flow weight (wavg.) & 66.24M & 83.92M & 49.40M & 51.67M & 181.14M & 206.45M & 136.46M & 180.42M \\
Out-flow weight HHI (wavg.) & 0.8311 & 0.0782 & 0.8082 & 0.0814 & 0.9496 & 0.0511 & 0.9397 & 0.0567 \\
\midrule
\multicolumn{9}{l}{\textit{Collateral financial features}} \\
Self-borrowing rate (wavg.) & 0.0645 & 0.0752 & 0.0455 & 0.0530 & 0.8661 & 0.2235 & 0.5350 & 0.3170 \\
Collateral utilization (wavg.) & 0.4117 & 0.0355 & 0.3800 & 0.0389 & 0.6424 & 0.0840 & 0.6630 & 0.0572 \\
Health ratio (wavg.) & 2.3947 & 0.3486 & 2.2618 & 0.3540 & 1.3535 & 0.2889 & 1.3780 & 0.2316 \\
Deposit share (wavg.) & 0.8240 & 0.0442 & 0.7716 & 0.0645 & 0.9005 & 0.0984 & 0.9008 & 0.0722 \\
\bottomrule
\end{tabular}

\end{threeparttable}
\end{table}


\begin{table}[ht]
\centering
\begin{threeparttable}
\caption{Descriptive statistics (Panel B) ---
         Borrower features,
         1 January 2021 -- 15 June 2024}
\label{tab:desc_stats_b}
\footnotesize
\begin{tabular}{lrrrrrrrr}
\toprule
& \multicolumn{2}{c}{ETH} & \multicolumn{2}{c}{WBTC}
& \multicolumn{2}{c}{DAI} & \multicolumn{2}{c}{USDC} \\
\cmidrule(lr){2-3}\cmidrule(lr){4-5}
\cmidrule(lr){6-7}\cmidrule(lr){8-9}
Variable & Mean & Std & Mean & Std & Mean & Std & Mean & Std \\
\midrule
\multicolumn{9}{l}{\textit{Borrower network features}} \\
In-degree (wavg.) & 3.6521 & 1.4271 & 3.7100 & 1.5268 & 2.3169 & 0.5319 & 2.2937 & 0.4527 \\
In-flow weight (wavg.) & 70.35M & 84.72M & 67.76M & 89.94M & 244.01M & 266.08M & 209.36M & 232.33M \\
In-flow weight HHI (wavg.) & 0.7682 & 0.1383 & 0.6725 & 0.1761 & 0.8268 & 0.0918 & 0.7856 & 0.0780 \\
Out-degree (wavg.) & 3.2595 & 0.9393 & 3.9186 & 1.1931 & 2.0375 & 0.3783 & 2.1007 & 0.3229 \\
Out-flow weight (wavg.) & 38.87M & 52.83M & 35.83M & 56.32M & 148.47M & 163.70M & 118.17M & 137.12M \\
Out-flow weight HHI (wavg.) & 0.7044 & 0.1192 & 0.6710 & 0.1493 & 0.8236 & 0.0833 & 0.7883 & 0.0747 \\
\midrule
\multicolumn{9}{l}{\textit{Borrower financial features}} \\
Self-borrowing rate (wavg.) & 0.4250 & 0.2251 & 0.4649 & 0.2557 & 0.4933 & 0.2169 & 0.3118 & 0.2024 \\
Collateral utilization (wavg.) & 0.5198 & 0.0771 & 0.5040 & 0.0850 & 0.5664 & 0.0684 & 0.5068 & 0.0657 \\
Health ratio (wavg.) & 1.6371 & 0.3213 & 1.6919 & 0.3209 & 1.6265 & 0.3156 & 1.7489 & 0.2562 \\
Deposit share (wavg.) & --- & --- & --- & --- & --- & --- & --- & --- \\
\bottomrule
\end{tabular}
\begin{tablenotes}[flushleft]
\setlength{\itemsep}{0pt}
\setlength{\parsep}{0pt}
\scriptsize
\item \textit{Notes}: Post-break sample, 1 January 2021 --
15 June 2024. Crypto columns use $s = 50\%$; stablecoin
columns use $s = 100\%$. Health threshold $= 10$: collateral
provider averages include only users with health ratio below
this value; results are insensitive to this choice within
the empirically relevant range. $\ell_{t,j}$ and $b_{t,j}$
are defined in Section~\ref{sec:empirical}; liquidity risk
is not reported for stablecoins as a full de-pegging
produces zero reserve losses by construction. Wavg.:
value-weighted average over collateral providers or
borrowers. HHI: Herfindahl--Hirschman index. Flow weights
in millions of dollars.
\item[$a$] Cross-pool variables are protocol-level and
identical across assets on any given date; reported once.
\end{tablenotes}
\end{threeparttable}
\end{table}

%\newpage
\FloatBarrier
\section{Simulated vs.\ Realized Liquidations}
\label{app:data_model}

\medskip

We compare simulated liquidation 
volumes from our propagation algorithm to realized 
daily liquidations on Compound over the full sample 
period. Realized liquidation volumes are obtained 
directly from the Compound V2 subgraph. For each day, 
we simulate liquidations using the observed end-of-day 
network snapshot under a range of shock sizes.

As discussed in Section~\ref{sec:Propagation_Algo}, our algorithm provides a conservative lower bound on realized liquidation volumes due to the close-factor discontinuity (Assumption~\ref{hyp:H1}) and the omission of intra-day liquidation cascades. To quantify this gap, we identify the price shock $s^* = s \times m$, where $s$ is the inter-day price shocks observed in the data, and $m \geq 1$ is a scalar multiplier that best matches realized liquidation volumes. 
%We select $m$ by maximizing the Pearson correlation between simulated and realized liquidation volumes over a grid $m \in [1.0, 4.0]$. As shown in \autoref{fig:correlation_simulated_realized}, the correlation reaches its maximum of $\approx 0.62$ within $m \in [1.75, 2.50]$, showing that realized liquidation volumes are best approximated when the inter-day effective price shocks are scaled up by $75$ to $150\%$, reflecting the amplification due to intra-day price fluctuations.
To assess whether it correctly identifies the days on which extreme liquidation activity occurs, we use the Precision-Recall AUC 
(PR-AUC) as our evaluation metric \citep{beutel2019machine, davis2006relationship, saito2015precision}.

We define a realized extreme day as one in the top $1\%$ of the realized liquidation-to-total-assets distribution, resulting in the selection of 17 days within our sample. For each 
multiplier $m \in [1.0, 4.0]$, we treat the simulated liquidation-to-total-assets ratio as a continuous risk score and evaluate it against the binary indicator of realized extreme days. The PR curve (\autoref{fig:pr_curve}) traces the tradeoff between two metrics as the flagging threshold varies. Precision measures the share of simulation-flagged days that are genuinely extreme in the data. Recall measures the share of the 17 extreme days correctly flagged. High thresholds produce few but reliable alerts (top-left of the curve), while low thresholds flag most extreme days at the cost of more false positives (bottom-right). A random classifier has PR-AUC equal to the base rate of $17/1{,}600 \approx 0.010$; a perfect classifier has PR-AUC 
$= 1.0$.

All multipliers substantially outperform the random baseline. The benchmark multiplier $m = 1.50$ achieves the highest PR-AUC of $0.309$ --- a $31\times$ improvement over random. At low recall thresholds (top-left of the curve), precision exceeds $60$--$80\%$: when the simulation flags only its most extreme days, the vast majority correspond to stress episodes in the data. This confirms that the simulation provides a reliable early warning signal for tail liquidation risk, even when calibrated using only end-of-day 
network snapshots and a single fixed shock size.

\begin{figure}[ht]
\centering
\captionsetup{justification=justified, singlelinecheck=false}
\includegraphics[width=.8\linewidth]{figures/precision_recall.png}
\caption{Precision-Recall curve --- simulation as early 
warning signal for extreme liquidation days.\\
\scriptsize{Note: The x-axis shows recall (share of 
the 17 realized extreme days correctly flagged); the 
y-axis shows precision (share of flagged days that 
are genuinely extreme). Each point corresponds to a 
specific threshold on the simulated 
liquidation-to-total-assets ratio. The dashed line 
is the random classifier baseline 
(PR-AUC $= 0.010$); $m = 1.50$ achieves the highest 
PR-AUC of $0.309$ ($31\times$ above baseline).}}
\label{fig:pr_curve}
\end{figure}

\autoref{fig:liquidations_stress} 
reports the six best-matched episodes, spanning the 
COVID crash (March 2020), the Bitcoin flash crash 
(May 2021), the January 2022 market correction, the 
Terra/Luna collapse (May--June 2022), and the FTX 
crisis (November 2022). In each case, the scaled 
simulated series closely tracks the shape and timing 
of realized liquidations, with the peak correctly 
identified on or near the day of the crash. This confirms that 
our propagation mechanism was economically active during 
realized stress events and that our simulations, once adjusted for 
intra-day fluctuations, provide a reasonable benchmark for the scale of 
losses. In accordance with this interpretation, the lower multiplier required for November 2022 
($m = 1.50$) relative to May 2021 ($m = 1.75$) is 
consistent with the more gradual nature of the 
FTX-induced price decline compared to the sharp 
Bitcoin flash crash.\footnote{Bitcoin's intraday high-low spread was approximately $30\%$ on May 19, 2021 and $15\%$ on November 9, 2022 
(source: CoinGecko historical OHLC data, \url{https://www.coingecko.com}).}
%Taken together, these results confirm that the propagation mechanism modeled here  was economically active during realized stress events and that our simulations provide a reliable benchmark for the scale and timing of losses that the network's structural features can generate.


%While full-sample correlation provides a useful baseline, our model is  designed to forecast systemic stress, making its performance during major  price crashes of particular relevance. We therefore also assess how well our model reproduces the large liquidation volumes observed in the wake of the two main crashes in our sample: the major sell-off of May 19, 2021 and the market crash of November 9, 2022 triggered by the FTX bankruptcy.  Applying episode-specific multipliers of $m = 1.75$ (May 19, 2021) and $m = 1.50$ (November 9, 2022) allows our model to closely match the realized liquidation volumes. \autoref{fig:liquidations_stress} reports the comparison for the two episodes. The simulated series closely tracks the shape and timing of realized liquidations in both cases, with the peak correctly identified on the day of the crash. 


\begin{figure}[ht]
\captionsetup{justification=justified, singlelinecheck=false}
\includegraphics[width=\linewidth]{figures/liquidations_stress_episodes.png}
\caption{Realized vs.\ simulated liquidations around six of the largest stress episodes in the sample.\\
\scriptsize{Note: Each panel shows daily liquidation 
volumes as a share of total protocol assets (\%) 
over a $\pm 7$-day window around the peak stress 
day (vertical dotted line). Dark blue solid: realized 
liquidations; orange dashed: simulated liquidations 
with episode-specific multiplier $m$ (shown in 
legend). Dollar values indicate peak realized 
liquidation volumes.}}
\label{fig:liquidations_stress}
\end{figure}

\section{Robustness Checks}
\label{app:Robustness_Checks}

\medskip

\subsection{Robustness to inclusion of borrower features ---
crypto assets.}
\label{app: Robustness_Borrower}

The baseline specification for crypto pools excludes borrower
characteristics from the candidate set, on the theoretical
grounds that risk transmission in crypto pools operates through
the collateral channel rather than through borrowers. When a
crypto price falls, it is the collateral posted by lenders that
becomes impaired, triggering liquidation cascades regardless of
borrower behaviour. To assess whether this exclusion materially
affects the results, we re-estimate the full pipeline for
$s = 50\%$ augmenting the candidate set with the borrower
network and financial features used in the stablecoin
specification.

\autoref{fig:borrowers_comparison} reports the comparison.
Adding borrower features leaves the collateral-side coefficients
largely stable in both sign and magnitude: no previously
significant variable changes sign, and the coefficient on
the dominant predictor --- value-weighted collateral utilization
($-0.0303$) --- is virtually unchanged. 

\FloatBarrier
\begin{figure}[ht]
\captionsetup{justification=justified,
              singlelinecheck=false}
\includegraphics[width=\linewidth]
    {figures/borrowers_crypto.png}
\caption{Robustness to inclusion of borrower features ---
trimmed model (no asset fixed effect), liquidity shock
$\ell_{t,j}$ (ETH and WBTC, $s = 50\%$).\\
\scriptsize{Note: Each column shows
the same three diagnostics for the two specifications side
by side. \textit{Row 1}: actual against fitted values on
the held-out test set. \textit{Row 2}:
HAC-estimated coefficients with 95\,\% confidence
intervals, ordered by magnitude; grey bars indicate
variables present in the union of both specifications but
not selected in that particular model. \textit{Row 3}:
residuals over time for the training (blue) and test
(coral) periods, separated by the dashed vertical line
at the train/test boundary (March 2023).}}
\label{fig:borrowers_comparison}
\end{figure}

\FloatBarrier
Out-of-sample predictive
performance is also unaffected, with test $R^2$ declining
marginally from 0.868 (without borrowers) to 0.827 (with
borrowers), consistent with the additional variables introducing
noise rather than signal in the crypto setting. LASSO retains
the same 14 collateral-side predictors in both specifications,
with the borrower variables either not selected or dropped at
the trimming stage. This confirms that borrower characteristics
carry no independent predictive content for crypto liquidity
risk once collateral-side features are controlled for, and
validates the parsimonious baseline specification.
\medskip


\subsection{Robustness to training subsampling.}
\label{app:Robustness_Subsampling}

\autoref{fig:robustness_size} reports the distribution of
out-of-sample $R^2$ across 50 subsampling draws for each shock
size $s \in \{10\%, \ldots, 90\%\}$, where each draw randomly
drops 10\,\% of training observations while keeping the test set
and variable set fixed. The figure therefore simultaneously
addresses two robustness questions: whether results are sensitive
to the choice of shock size, and whether they are sensitive to
which specific training observations are included.

\FloatBarrier 
\begin{figure}[ht]
\captionsetup{justification=justified,
              singlelinecheck=false}
\includegraphics[width=\linewidth]
    {figures/robustness_size_shock.png}
\caption{Robustness to random subsampling of the training set ---
trimmed$+$FE model, outcome: liquidity loss share $\ell_{t,j}$
(ETH and WBTC).\\
\scriptsize{Note: Each box plot summarises the distribution of
$R^2$ across 50 draws in which 10\,\% of training observations
are dropped at random; the red dot marks the baseline result
using the full training set. \textit{Left}: out-of-sample
(test) $R^2$. \textit{Right}: in-sample (training) $R^2$.
%For $s \leq 20\%$ model fit is low and variable because few
%positions are liquidated and the outcome has limited variation.
%From $s = 0.40$ onward, both train and test $R^2$ stabilise
%and the interquartile range collapses, confirming that results
%are insensitive to which specific training observations are
%included. The $s = 0.50$ benchmark, used for the main results,
%lies within this stable regime.
}}
\label{fig:robustness_size}
\end{figure}
\FloatBarrier 

On the first question, the figure confirms an inverted-U pattern
consistent with \autoref{fig:r2_by_shock_size}: model fit is
low and highly dispersed for $s \leq 20\%$, where few positions
are liquidated and the outcome has limited cross-sectional
variation. From $s = 40\%$ onward, the baseline test $R^2$
stabilises, with the $s = 50\%$ benchmark corresponding to the
peak of predictive accuracy.

On the second question, the box plots are narrow across all shock
sizes from $s = 40\%$ onward, with the baseline red dots sitting
on or near the median in every case. This confirms that the baseline results are not driven by specific 
training observations and that there is no indication of temporal 
parameter instability that would motivate a fixed-width rolling window 
over the expanding-window design.
The exception is $s = 10\%$, where the wider distribution simply
reflects the low-signal environment noted above.

Note that the figure is generated using the trimmed$+$FE
specification for comparability across shock sizes. Results are
virtually identical for the preferred no-FE specification, given
the non-significance of the fixed effect ($\hat{\gamma} =
-0.0039$, $p = 0.337$) and the negligible difference in test
$R^2$ (0.868 vs 0.871).


% ════════════════════════════════════════════════════════════════
\section{Model Comparison}

\medskip

\subsection{Liquidity Risk}
\label{app:model_comparison_liq}

\autoref{tab:model_comparison} summarizes the predictive
performance of all specifications considered for the liquidity
loss share $\ell_{t,j}$. The full OLS benchmark on all 30
predictors achieves the lowest training $R^2$ relative to its
complexity, with a test $R^2$ of 0.784 substantially below the
LASSO-based specifications, confirming that unrestricted OLS
overfits the training data. The LASSO model at $\lambda_{\min}$
selects 20 variables and achieves the highest test $R^2$ (0.878),
while the trimmed model retains only the 14 variables significant
at $\alpha = 0.001$ with a test $R^2$ of 0.868 and only a
marginal increase in RMSE. The asset fixed effect for WBTC is
not significant ($\hat{\gamma} = -0.0039$, $p = 0.337$) and its
inclusion in the trimmed model improves test $R^2$ by only 0.003,
confirming that the preferred specification is the trimmed model
without fixed effect.

\FloatBarrier 
\begin{table}[ht]
\centering
\begin{threeparttable}
\caption{Model comparison ---
         Liquidity Risk ($\ell_{t,j}$),
         ETH and WBTC, $s = 50\%$}
\label{tab:model_comparison}
\begin{tabular}{lrrrr}
\toprule
Model
    & Train $R^2$
    & Test $R^2$
    & RMSE
    & Vars \\
\midrule
Full OLS (all vars)
    & 0.855 & 0.784 & 0.0263 & 30 \\
LASSO $\lambda_{\min}$, no FE
    & 0.839 & 0.878 & 0.0198 & 20 \\
Trimmed, no FE (preferred)
    & 0.835 & 0.868 & 0.0205 & 14 \\
Trimmed $+$ FE
    & 0.836 & 0.871 & 0.0203 & 15 \\
\bottomrule
\end{tabular}
\begin{tablenotes}[flushleft]
\setlength{\itemsep}{0pt}
\setlength{\parsep}{0pt}
\scriptsize
\item \textit{Notes}: OLS on training set; Newey--West HAC
errors ($L = 4$ lags); test metrics on held-out test set.
WBTC fixed effect not significant ($\hat{\gamma} = -0.0039$,
$p = 0.337$); excluded from preferred specification.
Trimming removes predictors with $p \geq 0.001$.
\end{tablenotes}
\end{threeparttable}
\end{table}

\autoref{tab:coefs_comparison} and \autoref{fig:appendix_fe_comparison} report the coefficient
estimates for the trimmed model without and with the asset fixed
effect. Coefficients are remarkably stable across the two
specifications: in no case does the sign change, and the
magnitude changes are economically negligible, confirming that
the exclusion of the fixed effect does not distort the estimated
predictive associations.

\FloatBarrier 
\begin{table}[ht]
\centering
\begin{threeparttable}
\caption{Coefficient comparison ---
         Trimmed model without and with asset fixed effect,
         liquidity risk ($\ell_{t,j}$)}
\label{tab:coefs_comparison}
\footnotesize
\begin{tabular}{lrrrr}
\toprule
& \multicolumn{2}{c}{No FE}
& \multicolumn{2}{c}{With FE} \\
\cmidrule(lr){2-3}
\cmidrule(lr){4-5}
Variable
    & Coef. & $p$-value
    & Coef. & $p$-value \\
\midrule
\multicolumn{5}{l}{\textit{Own-pool characteristics}} \\
Collateral utilization
    & $+0.0131$ & $<0.001$
    & $+0.0130$ & $<0.001$ \\
Self-borrowing rate
    & $-0.0105$ & $<0.001$
    & $-0.0107$ & $<0.001$ \\
Utilization ratio
    & $+0.0057$ & $<0.001$
    & $+0.0060$ & $<0.001$ \\
Collateral factor
    & $-0.0122$ & $<0.001$
    & $-0.0105$ & $<0.001$ \\
\midrule
\multicolumn{5}{l}{\textit{Cross-pool linkages}} \\
In-degree dispersion (cross-pool)
    & $+0.0106$ & $<0.001$
    & $+0.0108$ & $<0.001$ \\
Out-degree (wavg.)
    & $-0.0083$ & $<0.001$
    & $-0.0083$ & $<0.001$ \\
\midrule
\multicolumn{5}{l}{\textit{Collateral network features}} \\
In-degree (wavg., collateral)
    & $+0.0074$ & $<0.001$
    & $+0.0073$ & $<0.001$ \\
In-flow weight (wavg., collateral)
    & $-0.0169$ & $<0.001$
    & $-0.0167$ & $<0.001$ \\
In-flow weight HHI (wavg., collateral)
    & $+0.0257$ & $<0.001$
    & $+0.0262$ & $<0.001$ \\
Out-degree (wavg., collateral)
    & $+0.0092$ & $<0.001$
    & $+0.0090$ & $<0.001$ \\
Out-flow weight HHI (wavg., collateral)
    & $+0.0094$ & $<0.001$
    & $+0.0093$ & $<0.001$ \\
\midrule
\multicolumn{5}{l}{\textit{Collateral financial features}} \\
Self-borrowing rate (wavg., collateral)
    & $-0.0043$ & $<0.001$
    & $-0.0040$ & $<0.001$ \\
Collateral utilization (wavg., collateral)
    & $+0.0303$ & $<0.001$
    & $+0.0300$ & $<0.001$ \\
Health ratio (wavg., collateral)
    & $+0.0153$ & $<0.001$
    & $+0.0151$ & $<0.001$ \\
\midrule
\multicolumn{5}{l}{\textit{Asset fixed effect}} \\
WBTC dummy ($d_a$)
    & \multicolumn{2}{c}{---}
    & $+0.0039$ & $0.337$ \\
\midrule
Constant
    & $+0.1037$ & $<0.001$
    & $+0.1018$ & $<0.001$ \\
\midrule
Train $R^2$
    & \multicolumn{2}{c}{0.835}
    & \multicolumn{2}{c}{0.836} \\
Test $R^2$
    & \multicolumn{2}{c}{0.868}
    & \multicolumn{2}{c}{0.871} \\
RMSE
    & \multicolumn{2}{c}{0.0205}
    & \multicolumn{2}{c}{0.0203} \\
\bottomrule
\end{tabular}
\begin{tablenotes}[flushleft]
\setlength{\itemsep}{0pt}
\setlength{\parsep}{0pt}
\scriptsize
\item \textit{Notes}: OLS on training set; Newey--West HAC
errors ($L = 4$ lags); predictors standardized using
training-set moments. Preferred specification: trimmed
model without WBTC fixed effect ($\hat{\gamma} = +0.0039$,
$p = 0.337$).
\end{tablenotes}
\end{threeparttable}
\end{table}

\FloatBarrier 
\begin{figure}[ht]
\captionsetup{justification=justified,
              singlelinecheck=false}
\includegraphics[width=\linewidth]
    {figures/liquidity_risk_predictions_comp.png}
\caption{Model comparison --- Trimmed model without and with
asset fixed effect, liquidity risk $\ell_{t,j}$ (ETH and WBTC,
$s = 50\%$).\\
\scriptsize{Note: Each row of panels shows the same four diagnostics
for the two specifications side by side.
\textit{Columns 1 and 3}: actual against fitted values on the
training set (blue) and the held-out test set (coral).
\textit{Columns 2 and 4}: HAC-estimated coefficients with
95\,\% confidence intervals, ordered by magnitude, and residuals
over time for the training (blue) and test (coral) periods
separated by the dashed vertical line at the train/test boundary
(October 2023).}
}
\label{fig:appendix_fe_comparison}
\end{figure}

% ════════════════════════════════════════════════════════════════
\subsection{Solvency Risk}
\label{app:model_comparison_solv}

\autoref{tab:model_comparison_solv} summarizes the predictive
performance of all specifications considered for the bad debt
share $s_{t,j}$. The full OLS benchmark on all 30 predictors
achieves a test $R^2$ of 0.835 but at the cost of using three
times as many variables as the preferred specification,
confirming that unrestricted OLS retains a large number of
redundant predictors without improving out-of-sample accuracy.
The LASSO model at $\lambda_{\min}$ selects 18 variables and
achieves a test $R^2$ of 0.822, while the trimmed model without
the fixed effect retains 13 variables with a test $R^2$ of
0.819. The addition of the asset fixed effect for WBTC improves
test $R^2$ substantially to 0.849 and reduces RMSE from 0.0076
to 0.0069, confirming that WBTC pools generate systematically
higher bad debt than ETH pools beyond what is explained by the
structural covariates. After post-HAC re-trimming at
$\alpha = 0.001$, the preferred specification retains 10
variables and is the trimmed model with the asset fixed effect.

\FloatBarrier
\begin{table}[ht]
\centering
\begin{threeparttable}
\caption{Model comparison ---
         Solvency risk ($b_{t,j}$),
         ETH and WBTC, $s = 50\%$}
\label{tab:model_comparison_solv}
\begin{tabular}{lrrrr}
\toprule
Model
    & Train $R^2$
    & Test $R^2$
    & RMSE
    & Vars \\
\midrule
Full OLS (all vars)
    & 0.867 & 0.835 & 0.0072 & 30 \\
LASSO $\lambda_{\min}$, no FE
    & 0.836 & 0.822 & 0.0075 & 18 \\
Trimmed, no FE
    & 0.834 & 0.819 & 0.0076 & 13 \\
Trimmed $+$ FE (preferred)
    & 0.838 & 0.834 & 0.0072 & 10 \\
\bottomrule
\end{tabular}
\begin{tablenotes}[flushleft]
\setlength{\itemsep}{0pt}
\setlength{\parsep}{0pt}
\scriptsize
\item \textit{Notes}: OLS on training set; Newey--West HAC
errors ($L = 4$ lags); test metrics on held-out test set.
WBTC fixed effect significant ($\hat{\gamma} = +0.0086$,
$p < 0.001$); retained in preferred specification.
Re-trimming after fixed effect inclusion reduces active
set from 13 to 10 predictors.
\end{tablenotes}
\end{threeparttable}
\end{table}
\FloatBarrier
\autoref{tab:coefs_comparison_solv} and \autoref{fig:appendix_fe_comparison_solv} report the coefficient
estimates for the trimmed model without and with the asset fixed
effect. The inclusion of the WBTC fixed effect produces
meaningful changes in several coefficients: the constant falls
from 0.0115 to 0.0071 as the asset-level intercept shift is
absorbed by the dummy, and the PageRank centrality coefficient
nearly doubles (from $+0.0033$ to $+0.0069$), reflecting that
the fixed effect disentangles the asset-level bad debt
differential from the network-position effect. Signs are stable
across all shared variables, confirming that the fixed effect
does not distort the estimated predictive associations but
sharpens their identification.

\FloatBarrier 
\begin{table}[ht]
\centering
\begin{threeparttable}
\caption{Coefficient comparison ---
         Trimmed model without and with asset fixed effect,
         solvency risk ($b_{t,j}$)}
\label{tab:coefs_comparison_solv}
\footnotesize
\begin{tabular}{lrrrr}
\toprule
& \multicolumn{2}{c}{No FE}
& \multicolumn{2}{c}{With FE (preferred)} \\
\cmidrule(lr){2-3}
\cmidrule(lr){4-5}
Variable
    & Coef. & $p$-value
    & Coef. & $p$-value \\
\midrule
\multicolumn{5}{l}{\textit{Own-pool characteristics}} \\
PageRank (scaled)
    & $+0.0033$ & $<0.001$
    & $+0.0069$ & $<0.001$ \\
Utilization ratio
    & $+0.0011$ & $<0.001$
    & \multicolumn{2}{c}{---} \\
In-weight HHI
    & $+0.0005$ & $<0.001$
    & \multicolumn{2}{c}{---} \\
\midrule
\multicolumn{5}{l}{\textit{Cross-pool linkages}} \\
Deposit concentration (HHI)
    & $+0.0013$ & $<0.001$
    & $+0.0015$ & $<0.001$ \\
Self-borrowing share (wavg.)
    & $-0.0013$ & $<0.001$
    & $-0.0015$ & $<0.001$ \\
Self-borrowing dispersion
    & $-0.0021$ & $<0.001$
    & \multicolumn{2}{c}{---} \\
Out-degree (wavg.)
    & $-0.0006$ & $0.003$
    & $-0.0010$ & $<0.001$ \\
\midrule
\multicolumn{5}{l}{\textit{Collateral network features}} \\
In-flow weight HHI (wavg., collateral)
    & $+0.0015$ & $<0.001$
    & $+0.0011$ & $<0.001$ \\
Out-flow weight (wavg., collateral)
    & $-0.0017$ & $<0.001$
    & $-0.0023$ & $<0.001$ \\
\midrule
\multicolumn{5}{l}{\textit{Collateral financial features}} \\
Self-borrowing rate (wavg., collateral)
    & $-0.0018$ & $<0.001$
    & $-0.0017$ & $<0.001$ \\
Collateral utilization (wavg., collateral)
    & $+0.0067$ & $<0.001$
    & $+0.0071$ & $<0.001$ \\
Health ratio (wavg., collateral)
    & $+0.0026$ & $<0.001$
    & $+0.0028$ & $<0.001$ \\
Deposit share (wavg., collateral)
    & $-0.0014$ & $<0.001$
    & \multicolumn{2}{c}{---} \\
\midrule
\multicolumn{5}{l}{\textit{Asset fixed effect}} \\
WBTC dummy ($d_a$)
    & \multicolumn{2}{c}{---}
    & $+0.0086$ & $<0.001$ \\
\midrule
Constant
    & $+0.0115$ & $<0.001$
    & $+0.0071$ & $<0.001$ \\
\midrule
Train $R^2$
    & \multicolumn{2}{c}{0.834}
    & \multicolumn{2}{c}{0.838} \\
Test $R^2$
    & \multicolumn{2}{c}{0.819}
    & \multicolumn{2}{c}{0.834} \\
RMSE
    & \multicolumn{2}{c}{0.0076}
    & \multicolumn{2}{c}{0.0072} \\
\bottomrule
\end{tabular}
\begin{tablenotes}[flushleft]
\setlength{\itemsep}{0pt}
\setlength{\parsep}{0pt}
\scriptsize
\item \textit{Notes}: OLS on training set; Newey--West HAC
errors ($L = 4$ lags); predictors standardized using
training-set moments. Preferred: trimmed model with WBTC
fixed effect ($\hat{\gamma} = +0.0086$, $p < 0.001$).
Variables marked~{---} excluded after post-HAC re-trimming
at $\alpha = 0.001$. Out-degree (wavg.) retained at
$p < 0.01$ in no-FE specification under relaxed threshold.
\end{tablenotes}
\end{threeparttable}
\end{table}

\FloatBarrier 
\begin{figure}[ht]
\captionsetup{justification=justified,
              singlelinecheck=false}
\includegraphics[width=\linewidth]
    {figures/solvency_risk_prediction_comp.png}
\caption{Model comparison --- Trimmed model without and with
asset fixed effect, solvency risk $b_{t,j}$ (ETH and WBTC,
$s = 50\%$).\\
\scriptsize{Each pair of columns shows the same four diagnostics
for one specification.
\textit{Columns 1--2}: Trimmed model without fixed effect.
\textit{Columns 3--4}: Trimmed model with WBTC fixed effect.
Within each pair: \textit{upper-left}: actual against fitted
values on the training set (blue); \textit{upper-right}: actual
against fitted values on the held-out test set (coral);
\textit{lower-left}: HAC-estimated coefficients with 95\,\%
confidence intervals, ordered by magnitude;
\textit{lower-right}: residuals over time for the training (blue)
and test (coral) periods, separated by the dashed vertical line
at the train/test boundary (March 2023).}}
\label{fig:appendix_fe_comparison_solv}
\end{figure}

% \begin{figure}[ht]
% \captionsetup{justification=justified,
%               singlelinecheck=false}
% \includegraphics[width=\linewidth]
%     {figures/appendix_borrowersT_comparison.png}
% \caption{Model comparison --- Trimmed model (no asset fixed
% effect) without and with borrowers' features, treasury loss
% share $\tau_{t,a}$ (ETH and WBTC, $s = 0.5$). Each column of
% panels shows the same three diagnostics for the two
% specifications side by side. \textit{Row 1}: actual against
% fitted values on the held-out test set (blue dots); the dashed
% line is the 45-degree line of perfect prediction. \textit{Row
% 2}: HAC-estimated coefficients with 95\,\% confidence
% intervals, ordered by magnitude; grey bars indicate variables
% present in the union of both specifications but not selected in
% that particular model. \textit{Row 3}: residuals over time for
% the training (blue) and test (coral) periods, separated by the
% dashed vertical line at the train/test boundary (October 2023).
% Adding borrowers' network and financial features leaves the
% collateral-side coefficients largely stable while providing
% additional explanatory power on the test set.}
% \label{fig:borrowers_comparison}
% \end{figure}

\newpage
% ════════════════════════════════════════════════════════════════
\section{\rev{Decision-Relevant Performance Metrics}}
\label{app:decision_metrics}

\rev{Test $R^2$ and top-decile precision, reported in
\autoref{tab:incremental}, measure average accuracy and the
quality of the induced ranking. This appendix reports three
further metrics that translate predictive performance into
quantities that matter for risk monitoring. All are computed on the
held-out test set of $401$ pool-days.}

\medskip

\rev{\boldpara{Magnitude of the tail error.} Because both outcomes
are expressed as a share of pool assets, a prediction error
translates into a dollar amount once multiplied by the assets of
the pool concerned. Restricting attention to the decile of
pool-days with the highest \emph{realized} stress, we report the
mean absolute error in millions of dollars per pool-day, together
with that error as a fraction of the loss actually realized on
those days.}

\rev{\boldpara{Alert budget.} Ranking all test pool-days by
predicted stress and working down the list, we record the share of
the panel that must be examined before $80\%$ of the realized
worst decile has been found. A perfect ranking requires examining
$8\%$ of the panel, since $32$ of $401$ pool-days is $8\%$.}

\rev{\boldpara{Tail calibration.} Within the realized worst decile
we report mean predicted stress divided by mean realized stress. A
value of one indicates that the model is unbiased where losses are
largest; a value below one indicates that severe outcomes are
systematically understated rather than merely predicted with
error.}

\FloatBarrier
\begin{table}[ht]
\centering
\begin{threeparttable}
\caption{\rev{Decision-relevant performance metrics.}}
\label{tab:decision_metrics}
\footnotesize
\begin{tabular}{lrrrrrr}
\toprule
& \multicolumn{3}{c}{\rev{Liquidity risk ($\ell_{t,j}$)}}
& \multicolumn{3}{c}{\rev{Solvency risk ($b_{t,j}$)}} \\
\cmidrule(lr){2-4}\cmidrule(lr){5-7}
\rev{Model}
    & \rev{MAE tail} & \rev{Budget $80\%$} & \rev{Calib.}
    & \rev{MAE tail} & \rev{Budget $80\%$} & \rev{Calib.} \\
\midrule
\rev{M1: Asset side}
    & \rev{\$36.5m ($40\%$)} & \rev{$57\%$} & \rev{$0.60$}
    & \rev{\$14.6m ($60\%$)} & \rev{$47\%$} & \rev{$0.40$} \\
\rev{M2: $+$ Pool features}
    & \rev{\$18.8m ($20\%$)} & \rev{$47\%$} & \rev{$0.81$}
    & \rev{\$8.9m ($37\%$)}  & \rev{$12\%$} & \rev{$0.64$} \\
\rev{M3: $+$ Collateral network}
    & \rev{\$10.2m ($11\%$)} & \rev{$25\%$} & \rev{$0.95$}
    & \rev{\$7.6m ($31\%$)}  & \rev{$12\%$} & \rev{$0.69$} \\
\rev{\textbf{M4: $+$ User features (full model)}}
    & \rev{\$13.5m ($15\%$)} & \rev{$24\%$} & \rev{$0.99$}
    & \rev{\$4.9m ($20\%$)}  & \rev{$11\%$} & \rev{$0.80$} \\
\bottomrule
\end{tabular}
\begin{tablenotes}[flushleft]
\setlength{\itemsep}{0pt}
\setlength{\parsep}{0pt}
{\scriptsize
\item \rev{\textit{Note}: Specifications M1--M4 are those of
\autoref{tab:incremental}. MAE tail: mean absolute prediction
error in millions of dollars per pool-day on the decile of
pool-days with the highest realized stress, with the same error
expressed as a share of the realized value in parentheses.
Budget $80\%$: share of test pool-days that must be flagged to
capture $80\%$ of that decile; the floor for a perfect ranking is
$8\%$. Calib.: mean predicted stress divided by mean realized
stress within the realized worst decile. Computed on the held-out
test set of $401$ pool-days, on which mean pool assets are
\$459 million.}
}
\end{tablenotes}
\end{threeparttable}
\end{table}
\FloatBarrier

\newpage
% ════════════════════════════════════════════════════════════════
\section{Predictive Performance Across Shock Sizes}
\label{app:shock_sizes}

\autoref{tab:model_comparison_allsizes} reports the predictive
performance of the trimmed model (no asset fixed effect) for
liquidity risk $\ell_{t,j}$ across the full range of shock sizes
$s \in \{10\%, 20\%, \ldots, 90\%\}$. The table confirms the
inverted-U pattern visible in \autoref{fig:r2_by_shock_size}:
test $R^2$ rises from near zero at $s = 10\%$ to a peak of 0.868
at $s = 50\%$, then declines as larger shocks push an increasing
share of positions into full liquidation, reducing the explanatory
power of the pre-shock network and financial characteristics. The
number of retained predictors also varies non-monotonically,
dropping to as few as 1 at $s = 10\%$ --- where the outcome has
near-zero variance --- and reaching 15 at $s = 90\%$, where a
small number of variables captures the residual variation among
the minority of positions not yet fully liquidated.

\FloatBarrier 
\begin{table}[ht]
\centering
\begin{threeparttable}
\caption{Predictive performance across shock sizes ---
         trimmed model (no FE), liquidity risk
         ($\ell_{t,j}$), ETH and WBTC}
\label{tab:model_comparison_allsizes}
\begin{tabular}{lrrrrr}
\toprule
$s$
    & $N$
    & Train $R^2$
    & Test $R^2$
    & RMSE
    & Vars \\
\midrule
0.10 & 2{,}008 & 0.079 & 0.116 & 0.0003 &  1 \\
0.20 & 2{,}008 & 0.460 & 0.463 & 0.0151 &  7 \\
0.30 & 2{,}008 & 0.666 & 0.661 & 0.0261 & 11 \\
0.40 & 2{,}008 & 0.816 & 0.799 & 0.0228 & 13 \\
\textbf{0.50} & \textbf{2{,}008} & \textbf{0.835} &
    \textbf{0.868} & \textbf{0.0205} & \textbf{14} \\
0.60 & 2{,}008 & 0.835 & 0.834 & 0.0215 & 14 \\
0.70 & 2{,}008 & 0.868 & 0.672 & 0.0247 &  7 \\
0.80 & 2{,}008 & 0.887 & 0.722 & 0.0107 &  9 \\
0.90 & 2{,}008 & 0.920 & 0.472 & 0.0055 & 15 \\
\bottomrule
\end{tabular}
\begin{tablenotes}[flushleft]
\setlength{\itemsep}{0pt}
\setlength{\parsep}{0pt}
\scriptsize
\item \textit{Notes}: Each row: separate model on training
set ($N_{\text{train}} = 1{,}607$), evaluated on held-out
test set ($N_{\text{test}} = 401$). LASSO at $\lambda_{\min}$
with HAC trimming at $\alpha = 0.001$ ($L = 4$ lags); no
WBTC fixed effect. Preferred specification ($s = 50\%$) in
bold. $s = 100\%$ omitted: no predictors survive trimming.
Low RMSE at $s \geq 0.80$ reflects near-zero outcome
variance, not improved accuracy; test $R^2$ confirms
deterioration.
\end{tablenotes}
\end{threeparttable}
\end{table}


% \section{Probit Model Fit}
% \label{app:probit_fit}

% \begin{figure}[ht]
% \captionsetup{justification=justified,
%               singlelinecheck=false}
% \includegraphics[width=\linewidth]
%     {figures/probit2.png}
% \caption{Probit model fit across shock sizes ---
% $P(\text{ETH hit harder than WBTC})$, own shock
% (blue solid) vs cross shock (coral dashed).
% \textit{Left}: ROC-AUC; the dotted horizontal line
% marks the random classifier baseline of 0.5.
% \textit{Right}: McFadden's Pseudo-$R^2$.
% Both metrics are consistently above their respective
% baselines across all shock sizes and both
% specifications, confirming that pre-shock pool
% characteristics provide meaningful discriminatory
% power for predicting which pool is hit harder.
% Own shocks are generally better predicted than
% cross shocks, consistent with the more direct
% transmission channel.}
% \label{fig:probit_fit}
% \end{figure}

% ════════════════════════════════════════════════════════════════
\section{Robustness to nonlinear specification}
\label{app:rf_benchmark}
% ════════════════════════════════════════════════════════════════

We estimate two random forest models for each outcome using the 
same chronological train/test split as the preferred OLS 
specification: one on the full candidate set of 30 variables 
and one restricted to the LASSO-selected variables. The 
comparison serves two purposes. First, it provides a model-free 
benchmark for predictive performance: if the linear OLS model 
outperforms the random forest out of sample, the linear 
specification is not missing important nonlinearities. Second, 
permutation importance scores from the full random forest 
provide an independent ranking of variable importance that does 
not depend on the linear specification.

\subsection{Liquidity Risk}
\label{app:rf_benchmark_liq}

\autoref{fig:rf_benchmark_liquidity} reports the results for 
liquidity risk. The trimmed OLS model achieves a test $R^2$ of 
$0.868$, outperforming both the full random forest ($0.773$\rev{\textsuperscript{?}}) 
and the restricted random forest ($0.820$). The superior 
out-of-sample performance of the linear model suggests that 
the relationship between the selected predictors and liquidity 
risk is well approximated by a linear specification, and that 
the random forest overfits the training data despite 
regularization. Permutation importance scores from the full 
random forest independently confirm the variable ranking 
implied by the OLS coefficients: collateral utilization of 
depositors and in-flow concentration emerge as the two dominant 
predictors, consistent with \autoref{fig:liquidity_predictions}, and 
variables excluded by LASSO rank uniformly lower. The agreement 
between the two approaches strengthens confidence that the 
LASSO-selected linear model identifies the true predictive 
structure of the data.

\FloatBarrier 
\begin{figure}[ht]
\centering
\captionsetup{justification=justified, singlelinecheck=false}
\includegraphics[width=0.8\linewidth]
    {figures/liquidity_RF.png}
\caption{OLS vs Random Forest benchmark ---
liquidity risk $\ell_{t,j}$ (ETH and WBTC, $s = 50\%$,
same chronological train/test split).\\
\scriptsize{Note:
\textit{Left panel}: permutation importance scores on the
held-out test set from a Random Forest estimated on the
full candidate set of 30 variables ($500$ trees,
\texttt{max\_features} $= \sqrt{p}$,
\texttt{min\_samples\_leaf} $= 5$).
Blue bars indicate variables selected by LASSO and retained
in the preferred OLS specification; grey bars indicate
variables excluded by LASSO.
Error bars show one standard deviation across 20 permutation
repeats.
\textit{Right panel}: test $R^2$ on the held-out test set
for three specifications --- trimmed OLS (preferred),
Random Forest on all 30 candidate variables, and Random
Forest restricted to the LASSO-selected variables.}}
\label{fig:rf_benchmark_liquidity}
\end{figure}

\subsection{Solvency Risk}
\label{app:rf_benchmark_solv}

\autoref{fig:rf_benchmark_solvency} reports the results for 
solvency risk. The trimmed OLS model again outperforms both 
random forests, achieving a test $R^2$ of $0.834$ against 
$0.727$\rev{\textsuperscript{?}} for the full random forest and $0.687$ for the 
restricted random forest. The larger gap between OLS and the 
restricted random forest --- compared to the liquidity case --- 
suggests that for solvency risk the predictive content of the 
LASSO-selected variables resides in their linear combination 
rather than in their individual marginal contributions, which 
the random forest exploits less efficiently.

The permutation importance ranking reveals an important nuance. 
The top-ranked variable by RF importance --- exposure-weighted 
collateral utilization (importance $= 0.117$) --- is also 
retained by LASSO, confirming it as the dominant predictor of 
solvency risk. However, several variables excluded by LASSO --- 
notably own-pool collateral utilization, In-degree, PageRank 
(raw), Out-degree, and the collateral factor --- also rank 
highly, unlike in the liquidity case where LASSO-excluded 
variables ranked uniformly low. The likely explanation is 
multicollinearity: LASSO dropped these variables not because 
they are uninformative, but because their predictive content 
overlaps with that of the retained variables, particularly 
the exposure-weighted collateral utilization and PageRank 
(scaled). The random forest, which selects splitting variables 
one at a time, recovers their marginal importance. The OLS 
still outperforms because the linear model already captures 
this shared information efficiently through the retained 
regressors. Notably, the WBTC fixed effect ranks tenth in 
permutation importance (importance $= 0.014$), confirming 
that it captures a genuine asset-level difference in solvency 
exposure.

\FloatBarrier 
\begin{figure}[ht]
\centering
\captionsetup{justification=justified, singlelinecheck=false}
\includegraphics[width=0.8\linewidth]
    {figures/solvency_RF_new.png}
\caption{OLS vs Random Forest benchmark ---
solvency risk $b_{t,j}$ (ETH and WBTC, $s = 50\%$,
same chronological train/test split).\\
\scriptsize{Note:
\textit{Left panel}: permutation importance scores on the
held-out test set from a Random Forest estimated on the
full candidate set of 30 variables ($500$ trees,
\texttt{max\_features} $= \sqrt{p}$,
\texttt{min\_samples\_leaf} $= 5$).
Blue bars indicate variables selected by LASSO and retained
in the preferred OLS specification; grey bars indicate
variables excluded by LASSO.
Error bars show one standard deviation across 20 permutation
repeats.
\textit{Right panel}: test $R^2$ on the held-out test set
for three specifications --- trimmed OLS (preferred),
Random Forest on all 30 candidate variables, and Random
Forest restricted to the LASSO-selected variables.}}
\label{fig:rf_benchmark_solvency}
\end{figure}



\clearpage
% Latest approved addition. Earlier Round 2 revisions retain their original colors.
\begin{latestrevisions}
\section{Institutional Comparison with Later Lending Protocols}
\label{app:later_protocols}

This appendix expands the comparison in Table~\ref{tab:protocol_designs}. It distinguishes three relationships: collateral that secures a borrower's debt, liquidity available to meet borrowing and withdrawals, and the accounts that bear unrecovered debt. A common user or collateral asset can expose several markets to the same shock without creating a contractual obligation for one market to absorb another's losses. The descriptions draw on official documentation accessed September 25, 2026. Historical features and later upgrades are identified separately.

\subsection{Compound III: one borrowable asset and protocol absorption}
\label{app:later_protocols_compound3}

Compound III organizes each deployment around one base asset that can be supplied or borrowed. Other accepted assets increase a user's borrowing capacity but are held as collateral rather than lent out. Suppliers of the base asset earn interest; supplied collateral does not earn interest from Compound III. Thus, a user can post several collateral assets against a base-asset debt, but cannot borrow a second asset from the same deployment.\footnote{\latestrev{Compound, \href{https://docs.compound.finance/collateral-and-borrowing/}{``Collateral \& Borrowing''}, especially ``Supply'' and ``Withdraw or Borrow.''}}

Liquidation also differs from V2. An outside participant calls the absorption function, which transfers the eligible account's debt and collateral to the protocol. The operation uses protocol reserves, and the protocol can subsequently sell the absorbed collateral for the base asset under its collateral-sale rules. This separates account liquidation from the later sale of collateral; it is not V2's exchange of debt repayment for cTokens followed by redemption.\footnote{\latestrev{Compound, \href{https://docs.compound.finance/liquidation/}{``Liquidation''}, including ``Absorb,'' ``Buy Collateral,'' and ``Reserves.''}}

For the present paper, the implication is that the collateral-pool withdrawal channel cannot simply be copied from V2: the collateral assets are not themselves borrowable pools within the deployment. The relevant stress test would track base-asset funding, the value of collateral backing that debt, and the absorption and sale process. Connections across separate deployments would need to be measured rather than inferred from a shared protocol name.

\subsection{Aave V3: shared reserves with version-specific restrictions}
\label{app:later_protocols_aave3}

An Aave V3 market contains multiple asset reserves. Within the permitted configuration, users supply assets, designate eligible collateral, and borrow other assets against their account. Liquidators repay eligible debt and receive collateral, with the amount governed by the deployed liquidation rules. This retains a user-level connection between collateral and borrowing reserves within a market. Separate market deployments have separate liquidity pools.\footnote{\latestrev{Aave, \href{https://aave.com/docs/aave-v3/concepts/liquidity-pool}{``LiquidityPool''}; \href{https://aave.com/docs/aave-v3/smart-contracts/pool}{``Pool''}, including the liquidation-call documentation.}}

The original V3 isolation mode limited a borrower to one isolated collateral asset, restricted borrowing to approved stablecoins, and applied a debt ceiling. It did not eliminate the connection between that collateral and the borrowed stablecoin reserves. Siloed borrowing was a different restriction: a user borrowing a designated siloed asset could not also borrow other assets. These controls therefore constrained the combinations and size of positions, rather than making every affected asset a financially independent lending market.\footnote{\latestrev{Aave, \href{https://aave.com/help/supplying/isolation-mode}{``Isolation Mode''}; Aave DAO's V3.7 documentation, \href{https://github.com/aave-dao/aave-v3-origin/blob/main/docs/3.7/mode-removal.md}{``Isolation Mode \& Siloed Borrowing Removal.''}}}

These descriptions must be dated. V3.7 removed the earlier isolation-mode and siloed-borrowing configurations, while introducing an isolated flag for efficiency-mode (eMode) categories. Those categories specify permitted collateral and borrowing combinations. Aave's changelog records the V3.7 rollout in April and May 2026. A current empirical application should therefore use the configuration and version in force for each market and observation date, rather than treat the original V3 controls as permanent features.\footnote{\latestrev{Aave DAO, \href{https://github.com/aave-dao/aave-v3-origin/blob/main/docs/3.7/Aave-v3.7-changelog.md}{``Aave V3.7 Changelog''}; Aave, \href{https://aave.com/docs/resources/changelog}{protocol changelog}, entries for April 20 and May 29, 2026.}}

\subsection{Aave V4: spoke-level positions and shared hub liquidity}
\label{app:later_protocols_aave4}

Aave V4 launched on Ethereum on March 30, 2026. Its hubs manage liquidity and accounting, while spokes manage user positions under their own configurations. A hub limits how much liquidity each spoke can supply or draw. This allows several sets of positions to use a common funding source without placing all users in a single collateral account.\footnote{\latestrev{Aave, \href{https://aave.com/docs/resources/changelog}{protocol changelog}, March 30, 2026; \href{https://aave.com/docs/aave-v4/liquidity/hubs}{``Hubs.''}}}

A user's health factor is evaluated within the relevant spoke. Positions held by that user in another spoke do not automatically support the first position or enter its health-factor calculation. Shared liquidity is therefore not equivalent to unrestricted cross-collateralization across a hub. A V4 stress test would need to distinguish the location of each user's collateral and debt from the hub through which liquidity is supplied and borrowed.\footnote{\latestrev{Aave, \href{https://aave.com/docs/aave-v4/positions}{``User Positions''}, especially ``Spoke Isolation.''}}

V4 liquidations seek a target health factor configured at the spoke level, rather than applying a single fixed close factor. The liquidation bonus varies with the position's health, and special rules govern small residual debt or collateral positions. This is related to the health-restoring approach in our simulation, but is not the same rule: our target is exactly one, our bonus is fixed in the exercise, and our proportional liquidation of multi-asset positions is an additional modeling assumption. Table~\ref{tab:liquidation_rule} consequently compares one dimension of liquidation design on the Compound V2 network; it does not estimate the overall performance of Aave V4.\footnote{\latestrev{Aave, \href{https://aave.com/docs/aave-v4/positions/liquidations}{``Liquidations''}, sections on the target health factor, variable bonus, and residual-position rules.}}

\subsection{Morpho Blue: isolated markets and multi-market vault exposures}
\label{app:later_protocols_morpho}

A Morpho Blue market pairs one collateral asset with one loan asset. Its parameters specify the collateral and loan tokens, price oracle, liquidation loan-to-value limit, and interest-rate model. Collateral secures borrowing in that market and is not lent out to its borrowers. Positions are assessed within the market rather than combined into a borrower's account spanning all Blue markets.\footnote{\latestrev{Morpho, \href{https://docs.morpho.org/learn/concepts/blue/}{``Variable Rate Market (Morpho Blue)''}; \href{https://docs.morpho.org/developers/contracts/blue/}{Blue contract documentation}, including the separate supply and collateral balances.}}

When a position exceeds the liquidation loan-to-value limit, a liquidator can repay part or all of its debt in exchange for collateral. The liquidation incentive depends on the market's specified limit. Any unrecovered debt is accounted for in that market and borne by its suppliers; it is not automatically charged to suppliers in a separate Blue market.\footnote{\latestrev{Morpho, \href{https://docs.morpho.org/learn/concepts/liquidation/}{``Liquidation on Morpho''}, including its discussion of bad debt.}}

Supply vaults create a separate layer of exposure. A vault can allocate a supplied loan asset across several underlying markets, so its depositors can be exposed to more than one collateral type. A loss in one allocation can affect the vault's depositors without causing another underlying market to inherit that loss. The timing and allocation of losses also depend on the vault version; the V1 documentation, for example, distinguishes how V1.0 and V1.1 recognize bad debt.\footnote{\latestrev{Morpho, \href{https://docs.morpho.org/learn/concepts/vault-v2/}{``Morpho Vault V2''}; \href{https://docs.morpho.org/curate/tutorials-v1/bad-debt/}{``Managing Bad Debt: How-to Guide''}, sections on Vault V1.0 and V1.1.}}

Applying the reconstruction approach to this design would therefore require market-level positions and, when vault depositors are the object of interest, the vault's allocation across markets. A common depositor or vault is evidence of shared exposure, not sufficient evidence of automatic contagion between otherwise isolated markets.

\subsection{What carries over from the Compound V2 analysis}
\label{app:later_protocols_implications}

The general contribution is a way to reconstruct obligations and assess which features of those obligations are informative about stress. The relevant balance sheets differ across designs. Compound III requires the absorption and reserve accounts; Aave V3 requires the deployed market's collateral and borrowing configuration; Aave V4 requires both spoke-level positions and hub-level funding; and Morpho Blue requires isolated market accounts, supplemented by vault allocations where relevant. These institutional differences guide a new application of the method. They do not support transferring Compound V2's coefficients, loss magnitudes, or ranking of predictive features to another protocol without estimation.
\end{latestrevisions}

\end{document}